\documentclass[12pt,openany]{book}

\usepackage{amsmath}      % 数学公式支持
\usepackage{booktabs}     % 三线表
\usepackage{tabularx}     % 自适应宽度表格
\usepackage{array}        % 表格列格式增强
\renewcommand{\tabularxcolumn}[1]{m{#1}} % 表格单元格垂直居中（上下居中）
\setlength{\extrarowheight}{2pt}          % 单元格上下边距对称补偿，确保视觉垂直居中
\usepackage{float}        % 固定浮动体位置 [H]
\usepackage{enumitem}     % 列表间距定制
\setlist[itemize]{
    itemsep=0.25\baselineskip,
    topsep=0.25\baselineskip,
    parsep=0pt,
    partopsep=0pt
}

% ========== 段内小标题定义（前后空0.5行） ==========
\newcommand{\subhead}[1]{%
    \par\vspace{0.5\baselineskip}%
    \noindent\textbf{#1}\par%
    \vspace{0.5\baselineskip}%
}
\usepackage{graphicx}     % 图片
\usepackage{multirow}

% ================= 导言区依赖 =================
\usepackage{tikz}
\usepackage{xcolor}
\usetikzlibrary{positioning, arrows.meta, calc}

% ================= 配色定义 =================
\definecolor{hqblue}{RGB}{28, 108, 168}
\definecolor{wqgreen}{RGB}{35, 142, 80}
\definecolor{cqorange}{RGB}{205, 90, 15}
\definecolor{bxred}{RGB}{185, 45, 40}

\definecolor{hqbg}{RGB}{244, 248, 253}
\definecolor{wqbg}{RGB}{243, 251, 246}
\definecolor{cqbg}{RGB}{254, 248, 242}
\definecolor{bxbg}{RGB}{254, 245, 245}

% ========== ETF 关系流程图配色 ==========
\definecolor{themeStockBg}{RGB}{242, 250, 244}
\definecolor{themeStockBorder}{RGB}{140, 200, 160}
\definecolor{themeStockText}{RGB}{35, 115, 65}

\definecolor{themeIndexBg}{RGB}{241, 246, 254}
\definecolor{themeIndexBorder}{RGB}{145, 185, 235}
\definecolor{themeIndexText}{RGB}{30, 90, 165}

\definecolor{themeFundBg}{RGB}{254, 248, 240}
\definecolor{themeFundBorder}{RGB}{240, 185, 130}
\definecolor{themeFundText}{RGB}{185, 90, 25}

\definecolor{lineColor}{RGB}{100, 115, 130}
\definecolor{outerBorder}{RGB}{210, 220, 230}
\definecolor{zoneInBg}{RGB}{252, 254, 255}
\definecolor{zoneOutBg}{RGB}{255, 253, 250}



% ========== 中文与字体（适配 Linux/Fedora + Fandol 字体族）==========
% 提示：Fedora 默认预装的 Noto CJK 为可变字体 (Variable Font TTC)，
% xdvipdfmx 暂不支持可变字体表结构，会触发 sfnt 错误导致编译中断。
% TeX Live 预装的 Fandol 开源字体族完整包含宋、黑、楷、仿宋，与 XeLaTeX 完美兼容。
\usepackage[fontset=fandol]{ctex}

% ========== 页面设置（适合阅读） ==========
\usepackage{geometry}
\geometry{
    paperwidth=145mm,      % 32开或A5感
    paperheight=210mm,
    left=20mm, right=20mm,
    top=25mm, bottom=25mm,
    headheight=15pt
}

% ========== 行距与段落 ==========
\usepackage{setspace}
\onehalfspacing                    % 1.5倍行距，阅读舒适
\setlength{\parindent}{2em}        % 段首缩进两字符
\setlength{\parskip}{0pt}          % 段间无额外间距
\usepackage{caption}
\captionsetup[table]{skip=0.5\baselineskip}
\raggedbottom                     % 避免垂直间距过度拉伸产生巨大空白
\setlength{\intextsep}{0.8\baselineskip plus 2pt minus 2pt}
\setlength{\textfloatsep}{0.8\baselineskip plus 2pt minus 2pt}
\renewcommand{\topfraction}{0.85}
\renewcommand{\bottomfraction}{0.85}
\renewcommand{\textfraction}{0.1}
\renewcommand{\floatpagefraction}{0.75}


% ========== 页眉页脚 ==========
\usepackage{fancyhdr}
\pagestyle{fancy}
\fancyhf{}
\fancyhead[LE,RO]{\thepage}
\fancyhead[RE]{\leftmark}
\fancyhead[LO]{\rightmark}
\renewcommand{\headrulewidth}{0.4pt}
\renewcommand{\chaptermark}[1]{\markboth{第\thechapter 章\quad #1}{}}

\fancypagestyle{plain}{
    \fancyhf{}
    \fancyfoot[C]{\thepage}
    \renewcommand{\headrulewidth}{0pt}
}

% ========== 章节样式（文学感） ==========
\usepackage{titlesec}
\titleformat{\chapter}[display]
    {\centering\Huge\bfseries}
    {第\thechapter 章}
    {20pt}
    {}
\titlespacing*{\chapter}{0pt}{-20pt}{40pt}

\titleformat{\section}
    {\Large\bfseries}
    {\thesection}{1em}{}

% ========== 修正 openany 下 mainmatter 强制 cleardoublepage 产生空白页问题 ==========
\makeatletter
\renewcommand\mainmatter{%
  \clearpage
  \@mainmattertrue
  \pagenumbering{arabic}%
}
\makeatother

% ========== 对话辅助 ==========
\newcommand{\dialogue}[2]{%
    \textit{——#1}\par
    \noindent\hspace{2em}#2\par\vspace{0.5em}}

% ========== 超链接与PDF书签 ==========
\usepackage[
    hidelinks,              % 隐藏超链接方框，保持版面清爽
    bookmarks=true,         % 生成PDF侧边大纲书签
    bookmarksnumbered=true  % 书签中保留章节序号
]{hyperref}

% ========== 文档信息 ==========
\title{\Huge\textbf{投资理财入门}}
\author{ToMercury \& Super Intelligence \\ \small\texttt{xmx2121@outlook.com}}
\date{2026年9月1日}

\begin{document}

% ========== 扉页与题献页 ==========
\frontmatter
\maketitle
\thispagestyle{empty}

% 题献页（独立单页，极简留白与居中美学）
\thispagestyle{empty}
\vspace*{\fill}
\begin{center}
    {\large\textit{—— \quad 献给自己的生日礼物}}
\end{center}
\vspace*{\fill}
\newpage

% ========== 免责声明与风险提示 ==========
\pagenumbering{Roman}
\setcounter{page}{1}

\chapter*{免责声明与风险提示}
\pdfbookmark[0]{免责声明与风险提示}{disclaimer}
\markboth{免责声明与风险提示}{免责声明与风险提示}

\small
\setlength{\parindent}{2em}

本书为作者个人学习投资理财过程中的知识归纳、实操复盘与经验总结，定位于财商教育通识科普，仅供学习、交流与参考。为保障读者的合法权益与资产安全，特作如下郑重声明：

\subhead{一、非投资建议声明}
本书所有内容（包括但不限于文字阐述、数据图表、计算公式、收益测算及推演案例）均不构成任何形式的投资建议、要约邀请、产品推荐或收益承诺。作者不具备证券投资顾问、基金投顾等法定金融从业资质，无意亦无权向读者提供任何实质性的投资决策指导。

\subhead{二、数据时效与客观局限}
书中引用的宏观经济数据、市场行情走势、交易规则、税费政策及监管制度均基于成书之时（2026年9月）的公开信息。由于资本市场与监管政策具有高度动态性，作者力求信息真实严谨，但不对内容的绝对准确性、完整性与最新时效性作任何明示或暗示的保证。任何实际操作均请以官方最新公告与法律法规为准。

\subhead{三、市场风险与盈亏自负}
市场有风险，投资需谨慎。任何金融资产均存在亏损本金的可能，过往业绩与历史统计绝对不代表未来表现。书中展示的复利测算模型仅为理论推演，绝不代表未来的实际确定回报。投资决策权完全属于读者个人，因阅读或参考本书内容进行的任何投资行为，其产生的一切风险与盈亏均由读者本人独立承担。

\subhead{四、中立立场与利益披露}
为便于读者理解相关机制与实操流程，书中提及了部分宽基指数、基金产品、基金公司、证券公司、第三方理财平台（如支付宝、天天基金、且慢等）及看盘交易软件。所有提及仅作为说明产品机制或行业生态的客观案例，不构成任何形式的背书或商业推广；作者与上述机构均无任何商业合作或利益分成关系，部分产品与平台仅涉及作者个人日常使用与实盘观察。

\subhead{五、专业咨询建议}
每个人的资产规模、风险承受能力、财务周期与现金流状况千差万别。若您面临具体且重大的财务规划或资产配置需求，请务必保持独立理性思考，并在必要时向正规持牌金融机构或具备资质的专业人士咨询。

\normalsize
\newpage

% ========== 目录 ==========
\begingroup
\setstretch{1.15}
\tableofcontents
\endgroup
\newpage

% ========== 前言/自序 ==========
\chapter*{前言}
\addcontentsline{toc}{chapter}{前言}
\markboth{前言}{前言}

国内教育体制中，很多基础教育面临严重缺失，比如“职业生涯教育”，“性教育”和“财商教育”。从某种程度来说，这些教育会潜移默化地影响一个人一生的轨迹，远远比一时的应试分数更为重要。可惜现实中，因为种种原因，很多人终其一生都未曾系统性接受过这些教育。而在这些缺失中，财商教育的空白尤为遗憾——因为我们一辈子都在和金钱打交道。

在我看来，“投资理财”理应是大学的通识课，人生的必修课。但是大部分父母不懂，学校又不教。我们大部分人按部就班地读小学、升中学、读大学，接着步入职场，赚钱、消费、结婚、贷款买车买房……整个教育与社会体系教会了我们如何出卖劳动力赚钱，也教会了我们如何消费花钱，却唯独没教过我们如何保护和管理自己的财富。正因如此，如何管理资产、建立正确的投资认知，便成了我们每个人自己的责任。

而正因为缺乏系统性的通识教育，社会大众对投资普遍存在极大的偏见与误解，往往简单地将投资等同于赌博，认为输赢全靠运气，甚至胜率早已被所谓的“庄家”暗中锁定。但事实恰恰相反，投资比拼的绝非赌徒手气，而是个人的认知框架、严守的纪律与成熟的心态。你不需要方方面面都做到极致，只需要在认知和心态上胜过随波逐流的普通人，就能在市场中稳健赚到属于自己的钱。

实际上，人的一生都在做各种不同的“投资”，只是很多人意识不到而已，比如买房买车、结婚生子。投资买房本质上就是借助贷款的“杠杆”，赌未来房价会继续上涨，从而赚取这部分收益。而杠杆又决定了风险和收益会被放大无数倍。然而，很多人对股市、投资避之不及，却对买房趋之若鹜，这何尝不是一种悲哀。

所以，与其说投资是一种赌博，不如说愚昧无知地盲目投资才是一种赌博。

我出生自河南某十八线小县城，家庭条件一般，父母对理财的认知更是为$0$，一辈子辛辛苦苦攒的钱全存在银行，吃着年化 $1.5\%$ 左右的利息，每年为银行经理送的冰箱彩电沾沾自喜，殊不知这点收益加起来远远跑不赢通货膨胀。

虽然我父母不懂理财，甚至自己舍不得花钱，但是他们对我花钱却很大方。大学期间，父母都是一次性给够一学期的生活费，通常还会有盈余。这也让我开始思考如何充分管理自己的钱财，同时我也有了足够的本金进行投资尝试。所以在这方面我非常感谢我的父母。

自从大学入学，我一直坚持学习投资理财。2022年10月刚入场时，市场还是熊市，彼时本金很少，没能在市场底部建立足够多的仓位。现在 2026年9月，市场早已进入了结构性牛市\footnote{结构性牛市：不同于全面牛市，结构性牛市是指股市中部分板块非常火热，冲到了历史高位，而其他板块十分冷清或有所下跌。例如近两年（2024年9月24日起）的科技板块非常火爆，但是其中泡沫相当大，随时有破裂的风险。}的中后期\footnote{2026年11月1日回顾：短短一两个月，熊市似乎已至。}，我也算经历了一轮简单的牛熊。在这轮牛熊之间，我一直坚持拿奖助学金和多余的生活费定投，但是奈何本金还是太少，所赚也并不太多。但是我觉得输赢并不重要，重要的是在投资中学习、复盘，提升自己的认知，建立自己的投资体系，这才是大学生投资的最大意义。

为什么写这本小书呢？
一来是市面上系统全面地讲解投资理财知识的内容实在太少，基本只聚焦于某一特定领域，比如只讲基金、股票或者ETF，但实际上它们之间的联系十分紧密，单独讲解就会非常割裂。
二来是本人大学四年所学的投资理财知识实在过于零散，写这本书也算是一种知识的整理归纳，构建自己系统性的知识体系。
最后是偶尔和同学交流投资理财，感觉大家对此一知半解，因此想要将这些知识分享给大家，希望大家一起学习思考，共同交流进步。

全书共分为五个章节，由浅入深，层层递进：

第一章“筑基篇”，剖析金钱的本质与真实通胀，揭示劳动收入与资本复利的巨大差距，搭建最基础的个人资产配置框架。

第二章“认知篇”，穿透企业成长与上市融资的逻辑，梳理一二级市场与各大板块架构，看清国家队、机构、游资与散户的博弈生态。

第三章“资产篇”，全景式拆解存款、债券、股票、基金与 ETF 等主流投资品，厘清底层逻辑与区别，并严正划清期货期权与金融骗局等危险红线。

第四章“实战篇”，立足普通人最切实际的时代红利，深入探讨全球宽基指数、美股跨境渠道与长期定投的实操心法。

第五章“心法篇”，结合作者自身的实操与观察，探讨本金底数、时间指数、消费陷阱乃至提前退休等深层次的人生与财富思考。

最后，这本小书因为篇幅和个人知识广度等原因，不可能涵盖投资的方方面面，只能作为最基础的入门参考。如果想要进一步学习投资理财知识，请自行查阅相关书籍或视频资料。

% ========== 正文 ==========
\mainmatter
\pagenumbering{arabic}
\setcounter{page}{1}

\chapter{筑基篇：财富底层逻辑}

正式步入理财之前，首先要建立正确的财商思维。当然认知的建立不是一蹴而就的，需要我们持续不断地学习。市面上有很多财商启蒙的书籍，比如《富爸爸穷爸爸》、《小狗钱钱》。当然大家也没必要看原书，直接用 AI 提炼核心观点学习即可。

下面是我个人的一些思考感悟，觉得比较基础和重要，分享给大家。

\section{金钱的本质与购买力}

\subhead{1. 破除“谈钱可耻论”}

货币的本质为一般等价物，就是我们出卖自己的劳动力换取货币，然后再购买他人的劳动力。货币只是中间产物，没有好坏之分。权力也是类似，权力只是实现某一目的的途径和手段。

经常在网上看到各种言论，比如“钱是万恶之源”、“可怕的权力”，其实真正可怕的不是权力和金钱，而是人性。

所以我们要破除“金钱可耻论”，摒弃“谈钱俗气”的想法，正确认识金钱，树立健康的财富观念。

\subhead{2. 认识“通货膨胀”}

通货膨胀会导致货币购买力下降，生活成本上升。国内通货膨胀率主要通过居民消费价格指数（CPI）的增长率来间接衡量。据统计，近 $20$ 年，CPI 年均增长率约 $2\%$。但是这和真实体感显然不符，因为官方CPI低估了真实购买力的贬值。官方CPI存在几个“不足”：
\begin{itemize}
    \item 房价不计入CPI（计入“房租”，而非购房成本）
    \item 资产价格（股票、房产）不计入
    \item 医疗、教育等价格涨幅远超CPI
\end{itemize}

相比之下，“M2增速-GDP增速”\footnote{M2增速-GDP增速：M2 是广义货币供应量，包括流通中的现金、企业活期存款、定期存款和居民储蓄存款等，常用来观察社会上的货币总量。该口径本身常被批评，这只是粗略的货币扩张参考指标，不能简单等同于通胀。另一个值得观察的指标为“GDP 平减指数”。}的统计体感更加真实。过去几十年，中国M2年均增速约 $10\sim 12\%$，GDP增速约 $6\sim 8\%$，两者差值约 $2\sim 6\%$。这是衡量货币超发、购买力稀释的更真实指标。

考虑到中国经济增速长期放缓，未来几十年的真实通胀率大概率前高后低，逐渐进入低通胀社会，类似日本。

所以银行里年化 $1\%$ 左右的收益，放在当前 $2\sim 6\%$ 的通胀面前，起不到任何作用。放的越多，亏的越多，放的越久，亏的越久。唯一的办法就是通过投资对抗通胀，甚至实现财富增值。

\section{劳动收入与资本复利}

\subhead{1. 线性劳动收入}

首先要明白一点，出卖劳动力所获得的收益是线性的，而资本增值的速度是指数型的。

人的一天只有 $24$ 小时，你再努力工作，也只会得到对应劳动时长的报酬（也可能得不到，比如自愿加班）。更何况人需要休息，工作时间有限，劳动报酬也就有限，甚至可能因失业、生病而中断。

相比之下，资本增值的速度远超常人想象，这背后的核心驱动力就是“复利”。

\subhead{2. 复利核心公式}

复利在数学上的表达非常简单，其计算公式为：

\[
F = P(1 + r)^t
\]

公式中只有四个变量，却揭示了财富积累的核心驱动机制：
$P$（Principal，期初本金）是复利运转的底仓与基石，年轻时本金较小是正常现象，关键在于尽早通过储蓄积累起第一笔启动资金；
$r$（Rate of return，年化收益率）是资产增值的加速器，但高收益往往伴随着高风险，巴菲特长达数十年的复合年化约为 $20\%$，成熟股市宽基指数的长期年化通常在 $8\%\sim 10\%$ 之间；
$t$（Time，投资期限）位于指数项上，是复利真正的“核动力”，只要给予足够长的时间跨度，哪怕初始本金不多，资产也能呈现出惊人的非线性增长；
$F$（Future Value，未来终值）则是我们最终积累的资产总额。

\subhead{3. $72$ 法则}

在日常投资中，如果想快速心算资产翻倍需要多少年，无需依赖复杂的对数运算，投资界有一个经典的估算工具——“$72$ 法则”：

\[
\text{资产翻倍所需年数} \approx \frac{72}{\text{年化收益率 } (r \times 100)}
\]

利用这个法则，我们可以快速看清不同资产增值效率的巨大差距：银行定期存款（年化约 $1\% \sim 1.5\%$）需要 $72 \div 1.5 = 48$ 年才能翻倍，很难跑赢真实通胀；稳健固收理财（年化约 $4\%$）需要 $72 \div 4 = 18$ 年翻倍；而核心宽基指数基金（长期年化约 $8\%\sim 10\%$）仅需 $7.2\sim 9$ 年即可实现翻倍。反过来，如果借了年化高达 $18\%\sim 24\%$ 的高息网贷或信用卡分期，欠款在 $3\sim 4$ 年内就会翻倍，这就是负债端复利的残酷之处。

\newpage

\subhead{4. 理想复利模型}

为了直观展现指数级增长的威力，下面给出两个简单的资本复利理想化数学模型。

第一个模型是“每月定投”：假设你每月月初投资 $1000$ 元（每年累计投入 $1.2$ 万元），有效年化收益固定为 $10\%$\footnote{有效年化$10\%$：主要股市长期收益参考，详见第四章。本模型按月复利，月收益率按 $(1+10\%)^{1/12}-1$ 折算，以保证有效年化收益率为 $10\%$。}，所有盈利分红全部用于再投资。在 $0\sim 40$ 年内，投入本金和总资产的变化曲线如图~\ref{fig:compound_curve} 与表~\ref{tab:compound_dca} 所示。

\begin{figure}[tbp]
    \centering
    \includegraphics[width=\textwidth]{../code/compound_interest/compound_interest_curve.pdf}
    \caption{定投复利曲线}
    \label{fig:compound_curve}
\end{figure}

\begin{table}[tbp]
    \centering
    \small
    \singlespacing
    \renewcommand{\arraystretch}{1.1}
    \setlength{\tabcolsep}{3pt}
    \caption{定投复利数据（每月投入1000元，有效年化10\%）}
    \label{tab:compound_dca}
    \begin{tabularx}{\textwidth}{|>{\centering\arraybackslash\hsize=0.6\hsize}X|>{\centering\arraybackslash\hsize=1.1\hsize}X|>{\centering\arraybackslash\hsize=1.1\hsize}X|>{\centering\arraybackslash\hsize=1.1\hsize}X|>{\centering\arraybackslash\hsize=1.1\hsize}X|}
        \hline
        \begin{tabular}{@{}c@{}}时间\\(年)\end{tabular} & 
        \begin{tabular}{@{}c@{}}总投入\\（万元）\end{tabular} & 
        \begin{tabular}{@{}c@{}}总资产\\（万元）\end{tabular} & 
        \begin{tabular}{@{}c@{}}总收益\\（万元）\end{tabular} & 
        \begin{tabular}{@{}c@{}}总收益/\\总资产\end{tabular} \\ \hline
        5  &  6.00 &   7.72 &   1.72 & $22.25\%$ \\ \hline
        10 & 12.00 &  20.15 &   8.15 & $40.43\%$ \\ \hline
        15 & 18.00 &  40.16 &  22.16 & $55.18\%$ \\ \hline
        20 & 24.00 &  72.40 &  48.40 & $66.85\%$ \\ \hline
        25 & 30.00 & 124.32 &  94.32 & $75.87\%$ \\ \hline
        30 & 36.00 & 207.93 & 171.93 & $82.69\%$ \\ \hline
        35 & 42.00 & 342.59 & 300.59 & $87.74\%$ \\ \hline
        40 & 48.00 & 559.46 & 511.46 & $91.42\%$ \\ \hline
    \end{tabularx}
\end{table}

由图表可以看到，第 15 年累计投资收益就已经超过了本金；到了第 40 年，累计收益占比更是超过了总资产的 $90\%$。这些数据看着似乎有些不可思议，但是事实就是如此，数学不会撒谎。

当然上述模型只是理想情况。首先，模型假设长期有效年化收益为$10\%$，且按月复利，而在现实中，市场是剧烈波动的，年化收益甚至可能为负。其次，模型并没有考虑通货膨胀对收益的侵蚀。最后，$40$ 年时间太久，长期预测没有任何意义。$40$ 年前苏联还没解体，谁又能预测 $40$ 年后的情况呢。当然，这个模型也能够说明一定的问题了。

第二个模型是“一次性大额投入”：假设手头有 $100$ 万现金，一次性全部投入。同样按照 $10\%$ 年化复合增长，那么第一年的收益就有 $10$ 万。看着可能不多，但是换算成 $8000$/月，据我保守估计，$80\%$ 的本科生达不到，$50\%$的研究生达不到。

如果本金 $100$ 万不变，每年的盈利分红全部用于再投资，具体收益见表~\ref{tab:compound_lump_sum}。

\begin{table}[tbp]
    \centering
    \small
    \singlespacing
    \renewcommand{\arraystretch}{1.1}
    \setlength{\tabcolsep}{3pt}
    \caption{一次性投入复利数据（本金100万元，年化10\%）}
    \label{tab:compound_lump_sum}
    \begin{tabularx}{\textwidth}{|>{\centering\arraybackslash\hsize=0.6\hsize}X|>{\centering\arraybackslash\hsize=1.1\hsize}X|>{\centering\arraybackslash\hsize=1.1\hsize}X|>{\centering\arraybackslash\hsize=1.1\hsize}X|>{\centering\arraybackslash\hsize=1.1\hsize}X|}
        \hline
        \begin{tabular}{@{}c@{}}时间\\(年)\end{tabular} & 
        \begin{tabular}{@{}c@{}}总投入\\（万元）\end{tabular} & 
        \begin{tabular}{@{}c@{}}总资产\\（万元）\end{tabular} & 
        \begin{tabular}{@{}c@{}}总收益\\（万元）\end{tabular} & 
        \begin{tabular}{@{}c@{}}总收益/\\总资产\end{tabular} \\ \hline
        5  & 100.00 &   161.05 &    61.05 & $37.91\%$ \\ \hline
        10 & 100.00 &   259.37 &   159.37 & $61.45\%$ \\ \hline
        15 & 100.00 &   417.72 &   317.72 & $76.06\%$ \\ \hline
        20 & 100.00 &   672.75 &   572.75 & $85.14\%$ \\ \hline
        25 & 100.00 & 1,083.47 &   983.47 & $90.77\%$ \\ \hline
        30 & 100.00 & 1,744.94 & 1,644.94 & $94.27\%$ \\ \hline
        35 & 100.00 & 2,810.24 & 2,710.24 & $96.44\%$ \\ \hline
        40 & 100.00 & 4,525.93 & 4,425.93 & $97.79\%$ \\ \hline
    \end{tabularx}
\end{table}

所以在我看来，“勤劳致富”和“富不过三代”不过是资本的谎言与穷人的自我慰藉，马太效应\footnote{马太效应：源自《圣经 $\cdot$ 马太福音》。强者能够利用已有财富、人脉等作为杠杆，撬动更大的资源，从而导致强者更强，弱者更弱。}在这里显现得淋漓尽致。

\section{负复利与税费常识}

然而，现实投资绝非表面上那么简单。很多人在理财中不仅没有享受到复利的红利，反而亏损累累。究其原因，是因为现实世界中存在着两个隐形的“复利杀手”——负复利与税费摩擦。

\subhead{1. 负复利}

复利既可以正向加速财富积累，也能在亏损时产生毁灭性的“负复利”。投资数学中存在一个极其残酷的规律：亏损与回本的涨跌幅是高度非对称的。

比如第一年本金暴涨 $50\%$，第二年暴跌 $50\%$，算术平均收益看起来是 $0\%$，但实际上本金经过计算（$100 \times 1.5 \times 0.5 = 75$）已经净亏损 $25\%$。

同理，如果本金亏损 $10\%$，需要上涨 $11.1\%$ 才能回本；亏损 $20\%$，需要上涨 $25\%$；亏损 $30\%$，需要上涨 $42.9\%$。而一旦亏损达到 $50\%$（资产腰斩），则需要暴涨 $100\%$（翻倍）才能回到原点；若不幸追高热门赛道或踩雷垃圾股导致亏损 $70\%$，需要暴涨 $233\%$ 才能回本；若亏损 $90\%$，则需要暴涨 $900\%$ 才能挽回损失。

这就是为什么巴菲特将“保住本金、控制回撤”列为投资的第一定律。大额亏损会直接斩断复利的链条，一次重大失误往往需要未来数年甚至数十年的平庸回报来弥补。

\subhead{2. 税务常识}

税收是国家财政收入的核心支柱，支撑着政府运转、国防安全、基础设施建设以及公共医疗教育等方方面面。但对于普通人而言，税收不仅是公民义务，更是贯穿我们一生的经济规则。

我国现行税制主要由流转税、所得税、财产税和行为税等大类构成。建立对税收体系的宏观认知，不仅能让我们看清商品价格背后的成本构成，也能读懂国家宏观政策的导向。国内常见的主要税种整理如表~\ref{tab:tax_system} 所示。

\begin{table}[tbp]
    \centering
    \footnotesize
    \singlespacing
    \renewcommand{\arraystretch}{0.92}
    \caption{我国主要常见税种一览}
    \label{tab:tax_system}
    \begin{tabularx}{\textwidth}{|>{\centering\arraybackslash\hsize=0.58\hsize}X|>{\centering\arraybackslash\hsize=0.88\hsize}X|>{\centering\arraybackslash\hsize=1.27\hsize}X|>{\centering\arraybackslash\hsize=1.27\hsize}X|}
        \hline
        类别 & 税种 & 主要征税对象/范围 & 特点 \\ \hline
        \multirow{2}{*}{流转税} & 增值税 & 销售货物、服务、无形资产、不动产及进口货物的增值额 & 第一大税种，环环征收、层层抵扣，最终由消费者隐性承担，一般 $6\% / 9\% / 13\%$ \\ \cline{2-4}
             & 消费税 & 烟、酒、高档化妆品、成品油、小汽车等特定消费品 & 引导消费、调节收入，在增值税基础上加征 \\ \hline
        \multirow{2}{*}{所得税} & \mbox{企业所得税} & 企业生产经营净所得（经纳税调整） & 基本税率 $25\%$（高新企业 $15\%$）\\ \cline{2-4}
             & \mbox{个人所得税} & 工资薪金、劳务报酬、股息分红与经营所得等 & 工资适用超额累进，投资所得享有专门政策 \\ \hline
        财产税 & 房产税 & 城镇经营性房屋（按余值或租金向产权所有人征收） & 个人非营业住房暂免（上海、重庆等试点除外）；从价 $1.2\%$ 或从租 $12\%$ \\ \hline
        \multirow{2}{*}{行为税} & 印花税 & 合同、产权转移书据、营业账簿、证券交易等应税凭证 & 征税面广、轻税重罚，日常经济活动中常见 \\ \cline{2-4}
             & 契税 & 土地房屋权属发生转移时，向承受方征收 & 由买方一次性缴纳 $3\%\sim 5\%$\textsuperscript{*}；置业刚性税负 \\ \hline
    \end{tabularx}
    \par\vspace{2pt}
    {\footnotesize\singlespacing\noindent\raggedright
    \textbf{注}：* 个人购买家庭唯一/第二套住房通常适用 $1\% \sim 2\%$ 的优惠税率（按面积和城市而定）。\par}
\end{table}

很多人常有“月薪三千不用纳税”的误解，这其实是混淆了“直接税”与“间接税”。工资未达起征点确实不用缴纳个人所得税（直接税），但只要发生消费，税收就如影随形。超市买米面粮油，价格内含增值税；买烟酒或汽车，除增值税，还叠加了高额消费税；买房置业，要一次性缴纳契税；如果未来房产税落地，每年还需要缴纳房产税。

我国税制结构以流转税为主，企业将税款内嵌在商品售价中，最终均由终端消费者隐性买单。理解了这一点，我们才能真正意识到，每个人无时无刻不在为公共财政贡献力量，金钱的管理与规划也绝不能脱离税收视角。

从宏观全貌回到投资理财，股票与基金的买卖，涉及印花税、个人所得税（如分红税）等。国家通过税制安排鲜明地鼓励“长期投资”，抑制“频繁投机”。其中涉及的具体规则，后文将进一步拆解。

此外，作为工薪族，要了解各种税收优惠政策。最直接的莫过于每年综合所得汇算清缴时的“专项附加扣除”，诸如 3 岁以下婴幼儿照护、子女教育、赡养老人、住房租金或房贷利息等，均可在税前按标准定额扣除，直接降低计税基数。另一项政策红利是“个人养老金”制度，每年可在 $12000$ 元限额内享受税前扣除，账户内投资收益暂不征税，退休领取时仅按 $3\%$ 单独计税。

合理合法地用好这些政策工具，相当于在没有任何市场风险的前提下直接省下税款，这也是工薪族增厚现金流与复利本金最具确定性的途径。

\section{个人资产配置框架}

搞清了复利的威力与现实中的挑战，我们就能明白为什么不能盲目梭哈。接下来在正式踏入市场之前，我们必须先为自己搭建起一套稳健、科学且攻守兼备的个人资产配置框架。

\subhead{1. 现金流管理}

在当今经济下行期，现金流非常重要。但是我发现身边的同学很难攒下自己的第一笔钱。可能是因为没有固定收入来源，也可能是因为消费欲望太高，入不敷出。

我觉得一方面我们可以利用手机 APP 记账，看清自己每个月的资金去向，减少非必要支出，为生活做减法。很多时候只有翻看账单，才能发现自己到底“浪费”了多少钱。

另一方面，我们可以先储蓄后消费。拿到生活费或工资，先把一部分强制存起来作为投资启动资金，剩下的作为当月可以支配的开销，以此积累人生的第一笔投资启动资金。

比如大学生可以每个月从生活费中抽出 $10\%$，大概 $100$ 多元，那么一年下来就能投入 $1500$ 元左右，虽然不多，但是用来学习刚刚好。即使一年亏掉 $10\%$，也只有 $100$ 多元，不会很肉疼。

\subhead{2. 资产配置“四笔钱”模型}

攒下启动本金之后，接下来的核心问题是：这些钱该如何分配？投资界有一句经典名言：“不要把所有鸡蛋放在同一个篮子里。”资产配置的本质，并非简单的分散购买，而是通过对资金属性的科学分流，在风险可控的前提下获取长期的复合回报。

我非常推崇“且慢 APP”\footnote{且慢 APP：盈米基金（中国证监会批准的独立基金销售机构）旗下一款投资交易软件。仅个人使用参考，不作为商业推荐。}所倡导的“四笔钱”资产配置理念：管好四笔钱，财富滚雪球。这四笔钱将我们的资金按照“时间期限”与“风险属性”划分为四个不同的账户，如图~\ref{fig:four_pots_matrix} 所示。

% ================= 四象限卡片代码 =================
\begin{figure}[tbp]
    \centering
    \begin{tikzpicture}[
        node distance=7pt and 8pt,
        card/.style 2 args={
            rectangle, 
            rounded corners=4pt, 
            draw=#1, 
            fill=#2,
            line width=0.7pt, 
            inner xsep=5pt, 
            inner ysep=6pt, 
            text width=4.5cm, 
            align=left,
            font=\linespread{1.15}\selectfont\footnotesize
        }
    ]
        % 活钱
        \node[card={hqblue}{hqbg}] (hq) {
            {\color{hqblue}\bfseries\small 活钱：随时要用的钱}\par
            \vspace{2pt}{\color{hqblue!40}\hrule height 0.5pt}\vspace{4pt}
            资金定位：每月开支、应急备用金\par
            投资品类：银行活期/存款、货币基金等\par
            核心特征：低风险低收益，流动性强
        };

        % 稳钱
        \node[card={wqgreen}{wqbg}, right=of hq] (wq) {
            {\color{wqgreen}\bfseries\small 稳钱：1--3年要用的钱}\par
            \vspace{2pt}{\color{wqgreen!40}\hrule height 0.5pt}\vspace{4pt}
            资金定位：买车款、年底旅游款等\par
            投资品类：债券、固收+、债券基金等\par
            核心特征：中风险中收益，流动性适中
        };

        % 长钱
        \node[card={cqorange}{cqbg}, below=of hq] (cq) {
            {\color{cqorange}\bfseries\small 长钱：3年以上不用的钱}\par
            \vspace{2pt}{\color{cqorange!40}\hrule height 0.5pt}\vspace{4pt}
            资金定位：买房资金、养老、子女教育等\par
            投资品类：股票、权益基金、指数ETF等\par
            核心特征：高风险高收益，需承受市场波动
        };

        % 保险
        \node[card={bxred}{bxbg}, below=of wq] (bx) {
            {\color{bxred}\bfseries\small 保险：防范极端风险的钱}\par
            \vspace{2pt}{\color{bxred!40}\hrule height 0.5pt}\vspace{4pt}
            资金定位：防范大病、意外等黑天鹅事件\par
            投资品类：社会保险（五险）、商业险（重疾险等）\par
            核心特征：以小博大，筑牢家庭安全网
        };
    \end{tikzpicture}
    \caption{资产配置“四笔钱”模型结构图}
    \label{fig:four_pots_matrix}
\end{figure}

很多散户在投资中之所以惨败，往往不是选错了标的，而是犯了“期限错配”的致命错误——拿随时要用的急钱去做了长期高风险投资。市场短期下跌之后，因为急用钱，不得不割肉离场。

四笔钱模型的巧妙之处，就在于通过清晰的边界划分，让每一笔钱各司其职。前三笔钱是我们在财富道路上冲锋的“矛”，而保险则是守护家庭后方的“盾”，二者缺一不可。

其中部分投资品类初学者可能不太熟悉，不用担心，后续章节我们会逐一深入拆解。现阶段，建立起“专款专用、攻守兼备”的四笔钱分流思维才是最核心的一步。

\subhead{3. 我的个人配置}

作为在校学生，我的日常开销相对固定，大部分钱短期内用不到。同时我的风险偏好较高，能承受较长时间的市场波动。

因此，我目前的资产配置比例大致为：$10\%$的活钱（年化1\%）、$30\%$的稳钱（年化5\%）、$60\%$的长钱（年化10\%），外加校园医疗保险。

按照资产加权计算，我的投资组合的综合年化收益为：$1\% \times 10\% + 5\% \times 30\% + 10\% \times 60\% = 7.6\%$。
这与我的实际收益基本吻合。

随着我近两年对市场的不断摸索学习，我发现我的投资组合仍然有很大的优化空间，所以理论上年化收益还能再往上提升。

特别需要提醒大家的是，市场上不存在放之四海皆准的黄金比例，只有最契合当下人生阶段的配置方案。随着年龄的增长，大部分人都会逐渐追求稳定、厌恶风险，这时便可以适当降低长钱的比重，降低资产波动。


\chapter{认知篇：资本市场运转机制}

上一章节中，我们了解到通货膨胀、资本复利和四笔钱模型等内容。我们明白了，想要实现资产增值，一定离不开债券、股票和基金等各类资产。但是上述所有资产的终极锚点只有一个——那就是“现代企业”。

所以本章节，我们将详细讲解“企业的诞生、发展和融资上市”等内容，然后进一步讲解“一级/二级市场”、“交易所体系”以及“机构与散户的博弈”，带领大家全面了解资本市场的运转机制。

\section{现代企业与上市}

\subsection{企业的成长路径}

\subhead{1. 企业融资}

假设有一家初创科技企业，想要研发人形机器人。众所周知，这类硬科技产品的研发周期长、资金消耗大，几位创始人凑齐的启动资金可能只够维持半年运转。此时企业摆脱资金困境通常有两种途径：一是向银行借贷，二是向外部资本寻求股权融资。对于缺乏抵押物、处于高风险研发期的初创公司，银行通常很难放贷，因此股权融资往往成为企业活下去并做大做强的必由之路。

融资的本质，是“用股权换取扩张资本”。假如最初始创始人有$100\%$的股权，公司发展的过程中，他拿出个人股份的$80\%$作为融资条件，那么最后上市，他就只剩下了$20\%$的股权。

很多人会疑惑：“创始人股权从$100\%$稀释到了$20\%$，这不是被资本‘割韭菜’了吗？”
其实不然。如果蛋糕只有 $100$ 万，拥有 $100\%$ 也只有 $100$ 万；但借助资本的杠杆把蛋糕做大到 $100$ 亿，即使只剩 $20\%$，身价也暴涨到了 $20$ 亿。现代商业的本质，就是学会“借力资本扩张，分散经营风险”。

\subhead{2. 融资上市流程}

初创公司的股权融资通常顺着企业生命周期，经历三个关键阶段：

第一阶段是天使轮/种子轮（Angel）。此时企业往往只有几个创始人、一份商业计划书（PPT）或一个粗糙的原型。死亡率极高，投资人主要看团队背景、领头人的个人魅力以及切入的赛道。

第二阶段是风险投资（VC），主要包括A/B/C轮融资。商业模式已被小范围验证，企业开始进入跑马圈地阶段。此时需要大量外部资本烧钱，搭建技术壁垒、抢占市场份额。

第三阶段是私募股权投资（PE），即Pre-IPO轮融资。企业已经进入成熟期，具备可观的营收与净利润。资本介入主要是为了规范财务、梳理合规风险，为上市做最后的冲刺。

\subsection{企业是否上市}

当企业发展壮大到成熟阶段，往往会面临一个关键抉择：首次公开募股（IPO，Initial Public Offering），即从小范围的非上市公司，变成面向全社会公开发行股票的上市公司\footnote{IPO上市：部分企业上市并不是为了做大做强，只是为了把包装精美的资产卖给二级市场，随后大股东高位减持套现离场，普通投资者需保持理性甄别。}。

上市被普遍视为企业发展的重要里程碑，但商业世界里也有华为、老干妈等企业始终坚持不上市。上市究竟能给企业带来什么好处？为什么有些企业却对其敬而远之？

\subhead{1. 上市的核心动因}

企业选择上市，主要能获得以下几项关键优势。

首先是获得无须刚性还本付息的权益资金。企业向银行贷款需要提供资产抵押，还要定期支付利息并到期还本。而股市 IPO 募集的是“权益资本”，只要企业持续合法经营，就无须向股民还本付息。这笔永久性资金为企业抵御经济周期与扩大再生产提供了极强的财务韧性。

其次是为早期投资人和创始团队提供变现退出渠道。非上市公司的股权流动性极差，很难在短时间内大额变现。一旦公司上市，股票在交易所每天都可以公开交易，陪伴企业多年的风投机构得以获利退出，创始团队的身价也能转化为具有流动性的真金白银。

此外，上市还能为企业带来显著的信用背书与并购便利。公众上市公司受到证券法和交易所的严格监管，定期披露审计财报，因此在银行贷款、政府招投标以及上下游供应链谈判账期时拥有更高的信誉度。

同时，上市公司的股票本身就是“流动资产”，企业后续可以通过“定向增发股份\footnote{定向增发：上市公司向符合条件的特定投资者非公开发行股份的融资行为，常用于并购同类资产或引入战略投资者。}”来收购竞争对手，不需消耗大量宝贵的现金流。

\subhead{2. 为什么有些企业坚决不上市}

上市虽然带来了充沛的资本，但同时也意味着企业要戴上多重制度枷锁。那些坚决不上市的企业，主要是在权衡上市所带来的副作用。

最主要的原因是避免被二级市场的“短期主义”绑架。上市公司必须按监管要求，每季度公布财务报告。然而二级市场机构往往极度看重当期利润，一旦某个季度业绩稍不及预期，股价就可能暴跌，甚至引发股东对公司日常经营横加干涉。以华为为例，作为一家全员持股的科技巨头，其核心技术研发动辄需要十年以上的战略定力和巨额资金投入，不上市正是为了避免被二级市场的短期盈利预期牵着鼻子走，确保长远战略不受资本干扰。

另一个重要前提是企业自身拥有强大的造血能力。例如老干妈坚持“不贷款、不欠款、现款现货”，日常现金流极其充沛，根本不需要依靠外部股市来融资。

此外，上市意味着公司的财务状况、核心客户和高管薪酬等商业信息全部公开暴露在聚光灯下，这对于处于激烈竞争中的某些行业而言，反倒会失去商业保密和战略灵活性的优势。

\section{证券市场基本架构}

明白了股票背后的企业实体，我们再来看承载这些股票交易的“舞台”——证券市场。

\subsection{一级市场与二级市场}

\subhead{1. 一级市场（发行市场）}

当企业经过多轮融资发展壮大，最终跨过上市门槛公开发行股票（IPO）时，机构与普通投资人开始认购公司新发行的股票，这部分发行承销的市场就是一级市场。

投资人付出的资金扣除承销费后，会实打实地注入上市公司的银行账户。普通散户除了参与“打新股”\footnote{打新中签：针对新上市的股票，可以以发行价进行申购，中签概率极低，本质上是一、二级市场的交接棒。一般情况下，首日开盘价会高于发行价，能够卖掉赚取打新收益。如果选择长期持有，新股破发（跌破发行价）并不罕见，需警惕风险。概念简单，但是细节较多，请自行了解。}，基本无法参与一级市场。

\subhead{2. 二级市场（流通市场）}

与一级市场企业的“公开发行”相对应，二级市场就是股票公开发行之后的“流通交易市场”，也就是我们日常在手机证券软件里看到的行情走势与交易盘口。

许多初学者常有一个误区：“我今天花了 1 万元买入某家上市公司的股票，是不是就等于向这家公司注资了 1 万元？”

事实完全不是这样。在二级市场中，绝大多数日常交易都只在投资者之间进行——比如买方看好某只股票，挂单买入，卖方正好想变现离场，交易系统撮合成交后，资金直接从买方的账户划转到卖方的口袋，上市公司一分钱都拿不到。因此，二级市场本质上是一个规范透明的“二手市场”，股票主要在机构与散户之间相互换手流转。

那么，二级市场上流通交易的股票究竟从何而来？自然源于一级市场发行后的股份。但这里存在一个初学者必须了解的实操常识：并非一家公司上市了，所有的股票都能在第一天上市流通。

在公司上市首日，股权通常分散在创始人、早期风投机构以及打新中签的散户手中。如果允许原始股东在上市当天就随便抛售，往往会引发大股东套现走人、股价踩踏崩盘的恶果。因此，监管机构对大股东与早期机构持有的股份通常设定了 $12\sim 36$ 个月不等的“锁定期”（即限售股）。在锁定期内，这些股票严禁在二级市场上抛售。

因此，上市初期在二级市场真正能够自由买卖的，主要是本次 IPO 中公开发行、面向市场的非限售流通股。直到锁定期结束（即“限售解禁”），大股东才能合规减持套现\footnote{减持/增持：股东看好自家公司发展，选择从二级市场收回自家股票，即为增持，一般视为积极信号。减持与之相反。}，此时通常会迎来二级市场的“解禁潮”，股价往往会出现大幅波动。

\subhead{3. 二者的紧密联系}

既然二级市场的钱无法流转到公司，那么设立二级市场的意义何在呢？

首先，二级市场是一级市场存在的“前提”，二级市场提供了一个退出机制。如果投资人获得的股票无法卖出变现，那么也就不会有风险投资，也就不存在一级市场了。

其次，二级市场给公司带来间接但真实的融资能力。公司可以通过定向增发、配股、发行可转债等方式进行再融资。

同时，公司股价可以作为抵押或信用，公司的股票可以拿来收购其他公司。所以二级市场是公司持续融资能力的定价基础和信用背书。

最后，二级市场还承担着投资风险转移、财富资源配置等重要作用。

\subsection{核心机构与监管体系}

\subhead{1. 证监会}

金融市场的核心监管机构。主要负责制定底线规则、打击内幕交易、查处财务造假、惩处市场操纵，充当市场的“警察与法官”。

\subhead{2. 证券交易所}

证券交易所是一个高度组织化、受严格监管的集中交易平台，相当于官方主办的“撮合集市”。

交易所自身不买卖股票，也不决定价格涨跌，其核心功能是制定标准化上市与交易规则、提供安全高效的自动撮合系统、公布实时行情数据，并对上市公司的日常信息披露和市场异常交易实施一线监管。

我们最常接触的 A 股、港股和美股主要交易所如表~\ref{tab:Exchanges} 所示：

\begin{table}[tbp]
\centering
\small
\singlespacing
\renewcommand{\arraystretch}{1.1}
\caption{主要交易所}
\label{tab:Exchanges}
\begin{tabularx}{\linewidth}{|>{\centering\arraybackslash\hsize=0.6\hsize}X|>{\centering\arraybackslash\hsize=1.6\hsize}X|>{\centering\arraybackslash\hsize=0.8\hsize}X|}
\hline
股市                  & 交易所       & 简称   \\ \hline
\multirow{3}{*}{A股} & 上海证券交易所   & 上交所  \\ \cline{2-3} 
                    & 深圳证券交易所   & 深交所  \\ \cline{2-3} 
                    & 北京证券交易所   & 北交所  \\ \hline
港股                  & 香港交易所     & 港交所  \\ \hline
\multirow{3}{*}{美股} & 纽约证券交易所   & 纽交所  \\ \cline{2-3} 
                    & 纳斯达克证券交易所 & 纳斯达克 \\ \cline{2-3} 
                    & 美国证券交易所   & 美交所  \\ \hline
\end{tabularx}
\end{table}

实际操作中，这三地的市场并非完全独立。许多中概股选择在美股与港股双重上市（如百度、阿里巴巴等），股票可以在两地交易并互相转换。

需要说明的是，境外上市架构通常分为“双重主要上市”（Dual Primary Listing）与“第二上市”（Secondary Listing）。前者两地均为第一主要上市地，须同时满足两地全部监管要求，其股票通常可直接纳入“港股通”\footnote{港股通：A股证券账户交易港股的渠道，账户资产大于50万元并通过相关考核即可开通。}以便境内资金交易；后者以另一市场为从属上市地，享有较多监管豁免，但在两地监管协调与股票互换机制上存在差异。例如阿里巴巴在 2024 年正式转为双重主要上市后，才得以正式纳入港股通。

\subhead{3. 券商与中国结算}

有了交易所和监管规则，普通投资者还不能直接进行交易，这背后还需要另外两个关键基础设施的支持。

一是证券公司（券商），充当投资者的交易代理通道。我们平时使用的手机炒股软件均由券商提供。投资者在券商开立证券账户，发出的买卖委托由券商系统高速报送至交易所撮合成交。为了防止券商挪用客户资金，我国严格实行“客户交易结算资金第三方存管制度”，投资者的交易本金由商业银行独立存管，券商仅能根据交易指令记账，无权直接提取资金。

二是中国证券登记结算公司（简称“中国结算”，CSDC），充当整个市场的“总记账员”。当交易所撮合成交后，究竟是谁买了谁的股票、资金如何划转清算，均由中国结算统一集中处理。这保证了无论我们在哪家券商开户，名下的股票资产权属都是绝对独立、安全且具有法律效力的。

\subsection{A股主要板块}

在了解了各大交易所之后，初入市场的投资者在看盘时往往会发现：不同股票的代码前缀大不相同（如 60、00、30、688 等），而且有的被称为“主板”，有的则叫“创业板”或“科创板”。

为什么同一个市场内部要划分出不同的板块呢？

这是因为不同企业在成长周期、经营体量、盈利能力与技术风险上天差地别。既有业绩稳健、现金流充沛的传统行业巨头，也有研发周期长、尚未盈利但成长空间巨大的前沿科技企业。如果用同一套上市门槛和交易规则去要求所有企业，要么会让大批代表未来生产力的新兴企业无法上市获得融资，要么会让普通散户在不知情的情况下承担过高的投资风险。

为此，我国依托不同交易所逐步建立了“多层次资本市场”，形成了定位互补的四大核心板块，其核心定位、交易权限门槛与涨跌停规则如表~\ref{tab:a_share_boards} 所示：

\begin{table}[tbp]
    \centering
    \footnotesize
    \singlespacing
    \renewcommand{\arraystretch}{1.08}
    \caption{A股主要板块核心特征与交易规则对比}
    \label{tab:a_share_boards}
    \begin{tabularx}{\textwidth}{|>{\centering\arraybackslash\hsize=0.85\hsize}X|>{\centering\arraybackslash\hsize=0.65\hsize}X|>{\centering\arraybackslash\hsize=0.65\hsize}X|>{\centering\arraybackslash\hsize=1.05\hsize}X|>{\arraybackslash\hsize=1.80\hsize}X|}
        \hline
        板块名称 & 股票代码 & 涨跌幅 & 准入门槛 & 核心定位与特色 \\ \hline
        沪深主板 & 沪60 / 深00 & $\pm 10\%$ & 无资产门槛 & 聚焦成熟期大型蓝筹\textsuperscript{1}与支柱国企，传统经济压舱石；代表方向为传统金融、能源与消费龙头。 \\ \hline
        创业板（深市） & 30开头 & $\pm 20\%$ & 10万资产 + 2年经验\textsuperscript{2} & 聚焦高新技术与高成长新兴产业，偏向创新创业与成长型实体；代表方向为新能源、高端制造等。 \\ \hline
        科创板（沪市） & 688开头 & $\pm 20\%$ & 50万资产 + 2年经验 & 面向世界前沿的“硬科技”，聚焦关键核心技术自主可控，允许未盈利上市；代表方向为集成电路、生物医药等。 \\ \hline
        北交所 & 8 / 920开头 & $\pm 30\%$ & 50万资产 + 2年经验 & 孵化创新型中小企业与细分赛道“专精特新”小巨人，企业体量较小，价格弹性与波动性最大。 \\ \hline
    \end{tabularx}
    \par\vspace{3pt}
    {\footnotesize\singlespacing\noindent\raggedright
    \textsuperscript{1} 蓝筹股：指经营业绩优良、现金流充沛、在所属行业中占有主导地位的大市值龙头企业。\\
    \textsuperscript{2} 准入门槛中的资金要求通常指申请开通权限前 20 个交易日证券账户内日均资产，交易经验自首次股票交易起算。各板块新股上市首 5 个交易日均不设涨跌幅限制。\par}
\end{table}

从表中可以看出，A股四大板块呈现出清晰的“风险与门槛匹配”特征：成熟稳健的主板不设资金门槛，涨跌幅限制在 $\pm 10\%$；而代表高成长、高研发风险的新兴产业与中小企业（创业板、科创板、北交所），涨跌幅限制逐步扩大到 $\pm 20\%\sim\pm 30\%$，同时设置了相应的资金验资与交易经验要求。这一制度设计，既为高成长科技创新企业提供了包容的融资土壤，也对风险承受能力有限的普通投资者构筑了防护栏。

关于 A股整体的交易机制（如 T+1 交易日、整手买入限制等）以及与美股等海外市场的深度横向对比，我们将在第三章“股票”小节中展开系统梳理（见表~\ref{tab:cn_us_stock_rules}）。

\section{市场参与者与核心角色}

在进入市场交易前，必须知己知彼，了解自己的对手。

\subsection{国家队的定位与职能}

散户口中的“国家队”（又称“汪汪队”、“狗队”），主要由中央汇金、中国证金、社保基金、外汇局旗下投资平台（梧桐树等）以及大型国有保险机构组成。

国家队的根本使命是平抑系统性金融风险，守护宏观金融稳定，同时实现国有资产的增值\footnote{国有资产增值：中国股市是一个政策市、融资市，国家队是要赚钱的，也要帮企业实现融资，那么钱从哪里来呢？}。

国家队通常会在市场遭遇系统性流动性危机、踩踏性恐慌抛售时“重拳出击”，买入巨额宽基ETF和超大型金融蓝筹股，向市场注入关键的流动性，阻断恐慌恶性循环。在市场情绪火爆，到达炒作高点时，减持或撤资，止盈离场。

所以股市中就有一种打法，专门选择追踪国家队动向，以此预估市场行情的高点和低点。
股市中追踪“国家队”的标的，主要分为国家队直接持有的宽基/行业ETF（如沪深300ETF华泰柏瑞），以及国家队重仓的个股（金融、银行、大型蓝筹股等）两大类。因为篇幅原因，具体细节不再展开，感兴趣的可以自行了解。


\subsection{机构、游资与散户}

当前二级市场的玩家主要分为下面三大派系。各路玩家的核心目标、最大优势与致命弱点可系统总结如表~\ref{tab:market_participants} 所示。

\subhead{1. 机构/主力}

机构主要包括公募基金、私募基金与外资等。以公募基金为代表的机构投资者，拥有强大的研究团队和海量资金，信奉基本面分析。

但他们有一个致命软肋——被动的基金份额赎回与严苛的相对排名考核。在牛市高点，基民\footnote{基民：并非股民，基民是指买基金的人。}贪婪追高涌入，公募被迫在最高位建仓接盘；在熊市底部，基民恐慌割肉赎回，公募被迫砸盘变现。这导致很多机构在客观上陷入了追涨杀跌的怪圈。


\subhead{2. 游资}

游资活跃在各大券商知名营业部，他们不看公司基本面，专精于题材炒作、情绪周期捕捉和封板战术。游资来去如风，赚的是市场博弈和击鼓传花的钱。如果散户跟着捕风捉影的消息去追高炒作，往往成为最后的接盘侠。

\subhead{3. 散户}

\begin{table}[tbp]
    \centering
    \small
    \singlespacing
    \renewcommand{\arraystretch}{1.15}
    \caption{证券市场各路玩家生态画像}
    \label{tab:market_participants}
    \begin{tabularx}{\textwidth}{|>{\centering\arraybackslash\hsize=0.8\hsize}X|>{\centering\arraybackslash\hsize=1.0\hsize}X|>{\centering\arraybackslash\hsize=1.0\hsize}X|>{\centering\arraybackslash\hsize=1.2\hsize}X|}
        \hline
        玩家角色 & 目标 & 最大优势 & 特点 \\ \hline
        国家队 & 防范系统风险、托底市场 & 资本实力雄厚、政策敏感 & 不以散户盈亏为导向 \\ \hline
        机构 & 赚取管理费、跑赢相对基准 & 投研实力强、合规规范 & 散户赎回压力、短期排名 \\ \hline
        游资 & 短线暴利、情绪博弈 & 速度极快、手法凶悍 & 严监管下易违规吃罚单 \\ \hline
        散户 & 赚快钱、盲目跟风 & 无强制平仓线、时间自由 & 信息严重滞后、追涨杀跌 \\ \hline
    \end{tabularx}
\end{table}

A股散户贡献了市场过半的交易量，却也是最容易亏损的群体。信息滞后、专业欠缺、本金分散、心态急躁，是散户常见的四大致命伤。

但是相比机构，散户也有巨大的优势。机构经理每年都要接受残酷的同行排名考核，私募有跌破净值的清盘红线，而普通散户既没有年终业绩考核，也没有被迫平仓的死线。

如果你使用的是第一章中配置的“闲置长钱”，你完全可以凭借“超长的时间容忍度”，在市场极度冰冷的周期低谷从容布局，耐心等待估值修复。以“耐心”对抗机构的“急躁”，以“周期”平复市场的“波动”，这是普通人在资本市场立于不败之地的真正底牌。

\chapter{资产篇：主流金融投资产品}

在上一章，我们大概了解了金融市场的运行逻辑，接下来我们讲解一下市面上常见的投资产品。

大部分学习资料，都习惯按“低、中、高风险”把投资产品机械分类。
但这种方式有一个致命的缺点：很多投资产品本质上是由更底层的资产组合或衍生而来的，单纯按风险高低机械分类，不仅容易割裂它们之间的内在联系，也容易让人只看表面标签，搞不清收益和风险到底源自何处。

本章节，我将从底层向上，一层层讲解不同的投资产品，让你对各个投资品类有系统全面的认识。

在讲解之前，我再次强调以下几点：

1. 风险和收益通常呈现正相关，想要追求更高的收益，就要承受更大的风险。市面上没有0风险产品，银行存款也不例外。

2. 本章节内容仅仅作为知识科普，介绍一些常见的基本概念，不构成任何投资建议！

\section{固收类资产}

固收类资产主要包括银行存款和债券。本质上都是把钱借给别人，然后收取利息。到期之后，你的收益就是“到期归还的本金+每年的利息”，基本是确定的，这也是固定收益（固收）这个名字的由来。固收类资产的特点是风险低、收益低。

\subsection{银行存款}

银行定期存款与大额存单是普通人最熟悉的储蓄与投资方式。以某银行 APP 中的数据为例，两种存款方式的特点大致如下：

\begin{itemize}
    \item 定期存款：起存金额低，取用相对灵活，当前利率一般为 $0.6 \sim 1.3\%$ 。
    \item 大额存单：一般 $20$ 万起步，利率略高于定期存款，一般为$1\% \sim 1.6\%$。
\end{itemize}

需要注意，上述存款方式都有固定期限，如果在到期前提前取走，一般按支取日活期利率计息。

我国自 2019 年起便已步入第四轮低利率周期。随着经济增速换挡与宏观利率下调，全社会的无风险收益率持续走低。在低利率周期下，单纯把资金死守在银行账户，微薄的利息收益很难跑赢真实的通胀与购买力贬值。

最后强调一点，大部分人觉得银行存款安全、0风险，但是实际上银行也存在破产或者跑路的风险（主要为小型村镇银行）。根据我国《存款保险条例》，当银行倒闭或破产时，本金和利息在 $50$ 万元以内可全额保障，超过部分按清算比例赔付。所以说，$50$ 万以内相对安全，$50$ 万以上可能存在风险。同时需要注意，银行的各类理财产品不属于存款保险条例的保障范围。

市面上从来没有0风险的产品，即使你仅仅追求 $1\%$ 的收益，也要承担一定的风险。

\subsection{债券}

\subhead{1. 什么是债券}

上一小节讲的银行存款，本质上是我们把钱借给银行，银行付给我们利息。而债券则把这个逻辑拓宽了——不仅仅是银行，国家、地方政府、金融机构、企业都可以向你“借钱”，然后给你一张“借条”，这张借条就是债券。

所以债券的本质就是一份标准化的债权凭证\footnote{标准化：意味着面值、利率、期限等条款都是统一格式的，可以在市场上流通交易，不像你借给朋友的钱那样“一事一议”。}。你买入债券，就是把钱借给了发行方；发行方按约定定期付你利息，到期归还本金。

债券和银行存款类似，有面值、利率和借债期限。国内债券面值通常为 $100$ 元/张\footnote{交易单位：实际交易中以“手”为单位，$1$ 手 = $10$ 张 = $1000$ 元面值。股票交易同样以“手”为单位，$1$ 手 = $100$ 股 。}。

\newpage

\subhead{2. 债券的分类}

债券按照发行主体的不同，可以划分为表~\ref{tab:bond_types} 中的几大类。

\begin{table}[tbp]
    \centering
    \small
    \singlespacing
    \renewcommand{\arraystretch}{1.25}
    \caption{债券分类}
    \label{tab:bond_types}
    \begin{tabularx}{\textwidth}{|>{\centering\arraybackslash\hsize=0.6\hsize}X|>{\centering\arraybackslash\hsize=1.3\hsize}X|>{\centering\arraybackslash\hsize=0.9\hsize}X|>{\centering\arraybackslash\hsize=1.2\hsize}X|}
        \hline
        类型 & 发行主体 & 风险等级 & 收益参考 \\ \hline
        国债 & 中央政府（财政部） & 极低 & $2\%$左右 \\ \hline
        地方债 & 省/市级地方政府 & 低 & 略高于国债 \\ \hline
        金融债 & 银行等金融机构 & 较低 & 略高于地方债 \\ \hline
        企业债 & 企业 & 中$\sim$高 & $3\%\sim 6\%+$ \\ \hline
    \end{tabularx}
\end{table}

简单来说，发行方的信用越高，违约风险越低，利率也就越低。国债背后是国家信用，除非国家破产，否则基本不会违约，所以利率最低。而企业债背后只有企业自身的经营能力做担保，企业经营不善就可能违约或跑路，所以需要用更高的利率来吸引投资者——这其实就是市场给风险定的“价格”。

此外，还有一种比较特殊的品种叫做可转债\footnote{可转债：全称“可转换公司债券”，持有人可以在特定条件下把债券转换成该公司的股票。它兼具债券的稳定性和股票的上涨空间，操作比较复杂，感兴趣可自行了解。}，它是债券和股票的“混血儿”，这里就不展开了。

\newpage

\subhead{3. 利率与债券价格}

前面说了，如果持有到期，收益基本是确定的。但如果你在到期之前想要卖掉呢？这时候债券的市场价格就会波动了。而影响债券价格最大的因素，就是市场利率。

这里有一个非常重要的规律：

\begin{center}
    市场利率上升 $\rightarrow$ 债券价格下跌

    市场利率下降 $\rightarrow$ 债券价格上涨
\end{center}

举个简单的例子。假如你去年花100元买了一张票面利率 $3\%$ 的债券（每年付3元利息）。今年央行加息了，市场上新发行的同类债券票面利率变成了 $5\%$（每年付5元利息）。

那么问题来了——现在市场上有人每年能拿5元利息，谁还愿意花100元买你这张每年只付3元的旧债券呢？所以你的旧债券价格就会下跌，可能只能卖95元甚至更低，直到买家觉得“以这个低价买入，持有到期的综合年化收益率和新债差不多”为止。

反过来也一样。如果央行降息，新债利率只有 $1\%$，你手里那张 $3\%$ 的旧债券就成了“香饽饽”，大家抢着买，价格自然水涨船高。

这就是“跷跷板”关系——利率和债券价格总是朝着相反的方向运动。

理解了这一点，我们就能搞明白一个很多人困惑的问题：债券基本稳赚不赔，那债券基金为什么会亏钱？

关键在于你买单只债券可以选择持有到期，中间的价格波动和你无关，到期拿确定的利息和本金就行。但是债券基金不一样——每天都有人买卖债券基金，基金净值是按照持仓债券的当日市场价格来折算的，不是按面值。

\subhead{4. 风险提示}

很多人觉得债券十分稳健，但实际上债券面临以下风险：

\begin{itemize}
    \item 信用风险：借钱的人还不起，违约了。国债基本不会违约，但企业债违约近年来时有发生，尤其是地产相关的信用债。
    \item 利率风险：市场利率上升会导致债券价格下跌。如果你急需用钱在利率上升期卖出，可能要亏本金。一般来说，期限越长的债券，利率风险越大。
    \item 流动性风险：部分债券交易不活跃，想卖的时候可能卖不出好价格，甚至短期内卖不出去。
\end{itemize}

固收不等于固赚，即使是追求稳健收益的“稳钱”，也要明白背后的风险逻辑，这样在市场波动时才不会恐慌操作、自乱阵脚。

\section{权益类资产}

\subsection{股票}

\subhead{1. 股票的本质}

上一节讲的债券，本质是你把钱借给别人，你是“债主”，赚的是利息。而股票则完全不同——当你买入一家公司的股票时，你不再是债主，而是成为了这家公司的股东。股票本质上就是一份标准化的企业所有权凭证。

既然成为了股东，你投资股票的收益主要来自以下两个方面：
\begin{itemize}
    \item 分红派息：公司经营赚了净利润，按持股比例把一部分现金打到你的账户里。
    \item 资本利得：股票市场中，低买高卖赚取的差价。
\end{itemize}

【税率补充】

红利税：分红需缴纳红利税，隶属个人所得税。上市公司现金分红对个人实行差异化税率，持股期限在 1 个月以内的，缴纳 $20\%$ 个税；持股 1 个月至 1 年的，缴纳 $10\%$ 个税；持股超过 1 年的，暂免征收红利税。

资本利得税：隶属个人所得税，目前我国对个人投资者在二级市场买卖股票和公募基金取得的差价收入，暂免征收个人所得税。美股对长期资本利得适用$0\%/15\%/20\%$累进税率，短期利得则按普通所得税率（最高37\%）征收。中国居民的美股交易差价所得，需在中国按“财产转让所得”申报，法定税率20\%。

印花税：A 股目前仅对股票出让方（卖出方）单边征收 $0.05\%$（万分之五）的印花税，买入方不征税；场内 ETF 买卖免征印花税。

\subhead{2. 除权除息}

很多初学者容易掉进一个认知陷阱：以为公司分红是白给的“红包”。事实上，股票在分红当天会进行除息（XD）。如果某股票价格为 $10$ 元，每股分红现金 $1$ 元，那么除息当天该股票开盘基准价会变成 $9$ 元，你的账户总资产在分红瞬间并没有任何变化。

分红不是天上掉馅饼，而是把公司资产直接折算成现金划归给你。长期来看，只有公司具备持续赚钱和增长的能力，分红和股价上涨才有真正的价值支撑。

与之类似的一个概念是除权（XR），公司向股东派发股票股利或进行送转股。例如公司每 $10$ 股送 $2$ 股，原本持有 $100$ 股的股东将变为 $120$ 股，同时股票价格按比例下降，原本每股 $36$ 元，现在就变为了 $36/1.2=30$ 元，公司总市值和个人总资产都不变。

既然除权除息把分红送股的价值从股价里扣掉了，股东账户总市值没变，那分红送股到底还有什么实际意义？
现金分红的意义在于提供一笔不被锁定的流动性。除息后股价降低，但股东手里多了可自由支配的现金，看好可低价再买入、持股数量增加，看空则已落袋、不必卖出变现。
送转股则降低了股价的名义门槛，拆细后能吸引更多中小投资者、改善流动性，但这本身不创造价值，脱离业绩的高送转也受监管限制。

归根结底，有意义的不是除权除息这个动作，而是此后公司能否用盈利增长把价格缺口“填”回来。持续赚钱，缺口会被利润逐步弥补，股东拿现金的同时股价也回到原位甚至更高；没有增长能力，除息后的低价只是对价值流失的确认，“填权填息”也就无从谈起。

\subhead{3. 价值与价格}

股票的买卖，都要经手交易所。需要强调的是，交易所是“撮合所”，不是“定价所”。
交易所只提供交易平台，供李四和王五相互买卖股票，至于股票的实时价格，完全是由供需关系和边际定价机制\footnote{边际定价：决定一家万亿市值公司当前“价格”的，往往只是此刻正在买卖的那几手股票的成交价。极少量的资金抛售，就足以将股价砸下数个百分点。详细可了解“交易所撮合机制”。}决定。

格雷厄姆曾有一句著名的投资名言：“股市短期来看是投票机，长期来看是称重机。”
在短期内，股票的“价格”波动由情绪、资金博弈和热点题材驱动，极度非理性；但在长期维度下，投机泡沫终将破裂，股价终究会向企业的真实商业盈利能力回归，体现出真实的“价值”。

因此，股市中可以简单划分出三大主流投资流派：

第一种是短线投资（趋势投资）。游资偏爱且散户最易跟风，核心依赖情绪博弈、题材炒作与技术指标，赚的是对手盘和“击鼓传花”的快钱。看似来钱快，但扣除交易成本后，本质是残酷的负和博弈，跟风散户极易成为最后的接盘侠。

第二种是长线投资（价值投资）。以巴菲特为代表，主张把买股票看作“入股企业”，赚的是公司长期利润增长与错误定价修复的钱。理论质朴但极度反人性，需要极高的商业认知与耐心。若盲目死扛业绩衰退的垃圾股，则可能沦为“价值毁灭”。

第三种是量化投资。以专业机构为主导，在美股等成熟市场中已占据大半成交量。它摒弃主观情绪，依靠数学模型、大数据回测与计算机毫秒级自动交易。在冷酷的算法与高频算力面前，散户靠肉眼和直觉做短线博弈无异于以卵击石。

以上三种投资方式并无优劣之分，关键在于投资者的认知、资金量、时间成本和风险承受能力。短线投机赚的是对手盘的钱，长线价值投资赚的是企业的利润增长和市场定价修复的钱，而量化投资赚的是算法和算力的钱。

\subhead{4. 市盈率与市净率}

既然长线价值投资看重企业的真实“价值”，那究竟该如何衡量一家公司的股票此刻是昂贵还是便宜？市场最基础的两把估值标尺是市盈率（PE）与市净率（PB）。

市盈率（$\text{PE} = \text{总市值} / \text{年净利润}$，即股价与每股收益之比），直观反映了按公司当前的盈利水平，买下它需要多少年能够收回成本。PE 越低通常代表资产越便宜。但需警惕的是，高成长公司往往享有更高的估值溢价，而强周期性行业在景气巅峰时因利润短期暴增，反而会呈现出极低的 PE，此时盲目投资极易踩入“低估值陷阱”。

市净率（$\text{PB} = \text{总市值} / \text{净资产}$，即股价与每股净资产之比），衡量的是股票价格相对公司账面净资产的溢价程度。它更适合银行、钢铁、重资产制造业等拥有大量固定资产且盈利容易受宏观波动影响的行业。

对于单只个股而言，估值受商业周期、竞争格局与财务调节的干扰极大，普通投资者很难单靠几个指标看清全貌。但这套估值逻辑一旦应用到由成百上千家核心企业组成的“宽基指数”\footnote{股票指数：简称指数，本质上是按照预设规则挑选出一篮子具有代表性的股票样本，并依据市值、股本或流动性等权重加权计算得出的综合价格指标（如同市场的“体温计”）。常见的核心宽基指数有沪深300、标普500等。跟踪指数买卖对应成分股的组合基金即为指数基金，第四章我们将系统展开讲解。}时，均值回归的规律就会展现出极高的实战价值，我们在第四章会展开详细运用。

\subhead{5. 股市交易规则}

股市中的交易规则非常繁琐复杂，后续碰到不懂的可以自行查阅，这里我们对比了解一下中美股交易机制即可，见表~\ref{tab:cn_us_stock_rules}。

注意，股市的交易机制一直在不断地修改完善，请以最新官方公告为准。

\begin{table}[tbp]
  \centering
  \small
  \singlespacing
  \renewcommand{\arraystretch}{0.92}
  \caption{中美股市核心交易规则对比}
  \label{tab:cn_us_stock_rules}
  \begin{tabularx}{\textwidth}{|>{\centering\arraybackslash\hsize=0.55\hsize}X|>{\centering\arraybackslash\hsize=1.2\hsize}X|>{\centering\arraybackslash\hsize=1.25\hsize}X|}
    \hline
    对比维度 & A股 & 美股 \\ \hline
    交易时间 & 工作日$9{:}30\sim 11{:}30$，$13{:}00\sim 15{:}00$；开盘集合竞价，盘后固定价格交易 & 北京时间 $21{:}30\sim 04{:}00$（冬令时延后$1$小时）；设盘前盘后交易 \\ \hline
    结算机制 & T+1 结算（当日卖出资金可继续买股票，次日方能取现） & T+1 结算（美股自 2024 年 5 月起全面实施 T+1） \\ \hline
    准入门槛 & 主板无门槛；创业板 $10$ 万+$2$年交易经验；科创板/北交所 $50$ 万+$2$年交易经验 & 一般无门槛 \\ \hline
    涨跌限制 & 沪深主板 $\pm10\%$，双创 $\pm20\%$，北交所 $\pm30\%$ & 个股无涨跌幅限制；大盘设有三级熔断机制 \\ \hline
    日内交易 & T+1，当天买入锁定，次日方可卖出 & 允许 T+0；小额账户受 PDT 规则\textsuperscript{1} 限制 \\ \hline
    交易单位 & 一般整手买入（至少 $1$ 手，即 $100$ 股）\textsuperscript{2} & 可买 $1$ 股，多数券商支持碎股（小数股） \\ \hline
    融资门槛 & 融资融券账户需 $50$ 万元资产门槛 & 融资账户通常只需 $2000$ 美元即可开通 \\ \hline
    主要风险 & 单向做多为主\textsuperscript{3}，T+1 限制下，买错无法当日纠错 & 无涨跌停 + T+0 + 高杠杆，单日亏损可能迅速放大甚至爆仓 \\ \hline
  \end{tabularx}
  \par\vspace{3pt}
  {\footnotesize\singlespacing\noindent\raggedright
  \textbf{注}：[1] PDT 规则：总资产不足 $2.5$ 万美元账户，5 个交易日内最多进行 3 次日内回转交易；[2] 科创板 200 股起购、北交所 100 股起购，均支持按 1 股递增申报，不足整手零股可一次性卖出；[3] A 股做空渠道有限，融券门槛高且券源稀缺，散户做空手段受限。\par}
\end{table}

\subhead{6. 致命红线}

在二级市场中，散户往往因为以下几大致命误区而亏光本金。

第一，盲信内幕消息，试图赚取认知之外的钱。网上大部分公开的“热心荐股”、内部代码，本质上都是利用散户流动性来高位出货或接盘。投资自己完全看不懂商业模式的企业，与在赌场押大小毫无分别。

第二，短线频繁交易与手续费磨损。许多散户误以为每天买进卖出赚取几个点是生财之道，但频繁交易面临三大致命缺陷：一是印花税与佣金等摩擦成本极高，一年下来可能吃掉 $10\%\sim 20\%$ 的本金；二是胜率极低，在随机波动的噪音中追涨杀跌会快速侵蚀本金；三是面对量化机构毫秒级的高频算法，散户的人肉交易在信息和反应速度上均被降维打击。

第三，盲目加杠杆炒股\footnote{杠杆：比如你本金10万元炒股，又贷款20万入场，那么现在就是3倍杠杆。如果股市跌10\%，你就亏损 $30\text{ 万元} \times 10\% = 3\text{ 万元}$，亏损本金的$30\%$。如果股价跌33\%，那就直接爆仓，本金全无。反之，盈利也是三倍。}。无论借钱、借网贷还是配资，杠杆虽然放大了收益，但更放大了波动与风险。一旦遭遇 $20\%\sim 30\%$ 的急跌，本金便会瞬间蒸发。在股市周期中，杠杆会让你在黎明到来之前，就被强制平仓踢出牌桌。

\section{组合类投资}

\subsection{基金}

\subhead{1. 什么是基金}

上一节我们讲了股票，直接买入单只个股不仅需要看财报、跟踪行业，还要时刻提防管理层造假和各种暴雷。对于普通散户而言，想要靠有限的资金去买齐几十只优质股票并做好风险对冲，难度极大。

而基金（通常指证券投资基金）的底层逻辑非常简单：凑钱办事、组合投资、专业打理、利益共享、风险共担。

简单来说，假如你手里只有 $2000$ 元，某些股票一手（$100$ 股）就要十几万元，你连一手都买不起，更不可能把资金分散到各行各业；但如果把一万个散户的 $2000$ 元汇聚起来，就形成了一个 $2000$ 万元的资金池。基金公司聘请专业的基金经理和研究团队，拿着这笔巨资按照既定策略去分散买入一整篮子资产（股票、债券、现金等）。你购买的基金份额，就代表你按比例拥有这一篮子资产的收益权。

\subhead{2. 基金投资的优点}

对于普通散户而言，基金之所以是极佳的投资工具，核心在于三大优势。

第一是低门槛分散风险。哪怕你只有 $10$ 块钱，也能通过基金间接持有一百家甚至几百家上市公司的资产。其中某一家公司暴雷亏损，对整个基金净值的影响不过是九牛一毛。

第二是专业省心。普通人有自己的本职工作，没有时间天天盯盘和分析枯燥的研报。把资金委托给专业团队管理，可以省去大量择股和频繁操作的心力损耗。

第三，资金安全与独立托管。公募基金实行严格的第三方托管机制。基民的资金并不存放在基金公司账户，而是存放在受监管的独立商业银行（托管行）。即使基金公司倒闭破产，基金资产也属于法定独立财产，清算后依然归全体持有人所有，不会被挪用抵债。

此外，公募基金有一套“严格”的信息披露制度。基金不仅每个交易日盘后公布当天的单位净值，还会定期披露财报。每季度结束后的 15 个工作日内披露季报，公布资产配置比例与前十大重仓股；每半年披露半年报，每年披露年报，公布全部持仓明细、隐形重仓股与费用开销。

\subhead{3. 基金的分类}

市面上的基金品类五花八门，我们可以通过表~\ref{tab:fund_classification} 快速理清基金的核心分类维度。对于普通投资者而言，日常打交道最多的是公募、被动指数型和宽基、红利类基金。

\begin{table}[tbp]
\centering
\singlespacing
\caption{证券投资基金核心分类}
\label{tab:fund_classification}
\footnotesize
\renewcommand{\arraystretch}{0.92}
\setlength{\tabcolsep}{3.5pt}
\begin{tabularx}{\linewidth}{|>{\centering\arraybackslash\hsize=0.55\hsize}X|>{\centering\arraybackslash\hsize=0.55\hsize}X|>{\centering\arraybackslash\hsize=1.9\hsize}X|}
\hline
维度 & 类型 & 核心特征与常见代表 \\ \hline
\multirow{2}{*}{募集方式} & 公募基金 & 面向大众（$1\sim10$ 元起投），强监管、信息公开（日常购买的主流） \\ \cline{2-3} 
                         & 私募基金 & 面向高净值合格投资人（通常百万元起），信息私密 \\ \hline
\multirow{4}{*}{持仓资产} & 货币型 & 只买存款、国债等，极低风险、高流动性（如零钱通、余额宝） \\ \cline{2-3} 
                         & 债券型 & $80\%$以上投向债券，稳健低波，适合追求稳健防守的资金 \\ \cline{2-3} 
                         & 股票型 & $80\%$以上投向股票，高收益与高波动并存，承担全部股市震荡 \\ \cline{2-3} 
                         & 混合型 & 股债仓位灵活配置（偏股混合/偏债混合/平衡），攻守兼备 \\ \hline
\multirow{2}{*}{管理方式} & 主动型 & 靠基金经理选股择时，力求跑赢大盘，管理费较高 \\ \cline{2-3} 
                         & 被动型 & 复制跟踪特定指数，也被称为指数基金\textsuperscript{1}，费率低且持仓完全透明 \\ \hline
\multirow{2}{*}{覆盖板块} & 宽基 & 不限单一行业，选样分散，代表整体经济或特定市值层级（如沪深300、标普500、中证A500） \\ \cline{2-3} 
                         & 窄基 & 集中押注单一特定行业赛道（如芯片、医药、白酒），受行业周期驱动、波动极大 \\ \hline
策略风格                 & 红利型 & 优选高股息、稳定分红的成熟防守企业（如银行、煤炭、公用事业），注重防御与现金流 \\ \hline
\multirow{2}{*}{投资地域} & 境内基金 & 投资国内A股及债券市场 \\ \cline{2-3} 
                         & QDII & 经特批投资海外市场（如纳指100、日经225），受外汇额度限制\textsuperscript{2} \\ \hline
\end{tabularx}
\par\vspace{2pt}
{\footnotesize\singlespacing\noindent\raggedright
\textbf{注}：[1] 指数基金包括宽基、窄基及策略指数（如红利）；[2] QDII 额度与外汇管制机制详见第四章。\par}
\end{table}

关于宽基指数与行业窄基的深层逻辑、长期表现差异以及为什么普通人应以宽基为主，我们在第四章第二节中有专项深入实战分析（见第~\pageref{sec:broad_based_funds}~页）。

\newpage

\subhead{4. 基金命名拆解}

截至 2026 年 9 月底，国内市面上共有 $14760$ 只公募基金\footnote{中国证券投资基金业协会，2026 年 9 月公募基金市场数据。}，数量甚至远超全市场的股票总和。面对这么多基金，最直接的筛选方法就是拆解它的命名格式：

\begin{center}
基金名称 = 基金公司 + 主题/风格 + 资产类型 + 份额类别
\end{center}

以“XX纳斯达克100ETF联接A”为例：

第一部分是基金公司，常见的头部公司有易方达、华夏、广发、南方、大成、汇添富、华泰柏瑞等。优先选择规模大、成立久的老牌公司，投研体系更健全。

第二部分是投资主题或跟踪标的，说明基金把钱投向哪里。比如“纳斯达克100”代表美股科技股，“沪深300”代表国内大盘，“医疗”、“半导体”则代表特定行业赛道。

第三部分是资产与运作类型，例如股票型、偏股混合型、债券型、ETF联接等。这一项直接决定了基金的底层股票仓位占比和预期波动大小。

第四部分是份额类别，通常以 A 类或 C 类呈现。同一只基金的 A 类和 C 类，底层持仓与基金经理完全相同，核心差异仅在于收费模式与持有周期。具体对比见表~\ref{tab:share_classes}。

\begin{table}[tbp]
\centering
\singlespacing
\caption{基金 A 类与 C 类份额对比}
\label{tab:share_classes}
\small
\renewcommand{\arraystretch}{1.2}
\begin{tabularx}{\linewidth}{|>{\centering\arraybackslash\hsize=0.7\hsize}X|>{\centering\arraybackslash\hsize=1.15\hsize}X|>{\centering\arraybackslash\hsize=1.15\hsize}X|}
\hline
对比维度 & A 类份额 & C 类份额 \\ \hline
申购费 & 买入时收取（平台通常打 $1$ 折） & 免费（0 申购费） \\ \hline
销售服务费 & 无 & 每日按净值计提（年化约 $0.2\%\sim 0.4\%$） \\ \hline
赎回费 & 持有越久费率越低，满 $1\sim 2$ 年免收 & 持有满 7 天通常即免收赎回费 \\ \hline
适合周期 & 长期持有（计划持有 1 年以上） & 短期持有（打算波段或持有数月） \\ \hline
总结口诀 & 拿得越久越划算，省下日常服务费 & 进出灵活成本低，省下买入申购费 \\ \hline
\end{tabularx}
\end{table}

\subhead{5. 风险提示与行业内幕}

我们把钱交给专业人士打理，是不是可以完全放心了？绝对不是！基民和基金公司存在天然的利益错位，我们来分析一下双方情况。

首先，我们主动把钱交给基金公司，基金公司赚取的是我们的手续费。无论我们的账户赚钱还是亏钱，基金公司每天都在旱涝保收地提走管理费。

其次，基金公司手握散户几十亿资产，每季度仅披露前十大持仓股，每半年披露一次全部持仓，操作空间较大。

最后，基金公司、基金经理与各个公司的股东、庄家可能存在千丝万缕的联系。

综合上述几点，主动基金可能存在以下风险\footnote{这些风险为中国证监会披露、网上流传或本人亲身经历。仅代表个人观点，不构成任何投资意见。}：

第一，老鼠仓风险（法律禁止）。金融从业人员利用内幕信息和散户的钱财进行的非法利益输送。

第二，基金风格漂移。风格漂移的原因有很多，可能是更换基金经理，投资理念不同导致漂移；可能是由于业绩压力，追逐热门行业导致风格漂移。

还有一种常见的，因为某些原因\footnote{原因可能有很多，比如内幕交易，或者QDII额度有限，一份美股 + 九份A股，就可以用一份QDII额度，赚取10份手续费。}，主观刻意而为的。比如基金声称主要投资海外资产，以此来吸引资金。但实则可能大部分都买了大A，侵蚀了投资者的个人收益。又比如某“红利型基金”，披露持仓基本全是科技股，许多不懂的小白就可能被套在其中。

第三，高管理费。主动基金的管理费通常高于被动基金，因为主动基金需要基金经理频繁管理和调仓。被动基金只需要被动追踪某一指数，管理费用较低。

第四，窄基行业集中、周期明显。比如买入了某“科技型”主动基金，在“科技牛市”中可能大涨。但基金风格固定、行业集中，基金经理即便预判科技股将要大跌，也受合同约束难以转投其他低估值板块。因此，窄基的择时责任更多落在投资者自己身上，需要关注市场变化，自行判断减仓或清仓的时机。

所以如果买主动基金，最好买正规公司、无黑历史、成立时间久、运作稳定的优质基金。当然，最推荐普通人或小白购买风险更小、更加省心的被动基金。

\newpage

【事实胜于雄辩】

\begin{figure}[tbp]
    \centering
    \includegraphics[width=0.78\textwidth]{../figure/category.png}
    \caption{某“高股息主题”基金行业分布披露}
    \label{fig:high_dividend_drift}
\end{figure}

众所周知，红利型基金/高股息基金，主要购买现金流稳定、股息率较高的成熟企业股票，如银行、金融、铁路、消费等老登股。这种基金通常波动较小，以稳健分红著称。

打开支付宝，搜索某只“高股息主题股票”基金，映入眼帘的近一年收益率达到了惊人的$-17\%$。一只主打高股息防守的基金，为什么会在市场震荡中跌成这样？

打开基金持仓，查阅 2026 年 6 月 30 日的行业分布披露，真相瞬间浮出水面。

如图~\ref{fig:high_dividend_drift} 所示，在这只名为“高股息”的基金中，
排名第一的行业竟然是“科技”，占比高达惊人的 $51.46\%$，占据了半壁江山。
随后是制造（$19.98\%$）与资源（$19.64\%$）。
而真正意义上的高股息传统主力——“金融”与“消费”，持仓占比居然全部为 $0.00\%$！请问此时监管何在？

\subsection{ETF}

\subhead{1. 场内与场外}

在深入 ETF 之前，首先必须搞清楚投资圈最常挂在嘴边的两个词——“场内”与“场外”。

这里的“场”，指的就是证券交易所（如上交所、深交所等）。

简单来说，我们平时在支付宝、微信或银行买基金，属于“场外交易”——本质是直接向基金公司申购或赎回，按天核算净值\footnote{补充：以交易日 15:00 为界，之前下单按当日收盘净值成交，之后顺延至次日，QDII 基金受海外市场时差与结算安排影响，确认和到账可能再延迟；部分平台截止时间可能更早，具体以销售平台公告为准。}；而在券商交易软件里像买卖股票一样与其他投资者实时竞价撮合，则属于“场内交易”。两者的核心机制与规则差异见表~\ref{tab:in_out_market}。

\begin{table}[tbp]
\centering
\small
\singlespacing
\caption{场内交易 vs 场外交易核心差异}
\label{tab:in_out_market}
\renewcommand{\arraystretch}{1.1}
\setlength{\tabcolsep}{5pt}
\begin{tabularx}{\linewidth}{|>{\centering\arraybackslash\hsize=0.55\hsize}X|>{\centering\arraybackslash\hsize=1.15\hsize}X|>{\centering\arraybackslash\hsize=1.3\hsize}X|}
\hline
对比维度 & 场内交易（二级市场） & 场外交易（开放式渠道） \\ \hline
交易场所 & 证券交易所（券商 APP） & 支付宝、微信、天天基金、银行等 \\ \hline
开户与否 & 需要开立证券账户 & 无需证券账户，绑定银行卡即可 \\ \hline
交易标的 & 股票、场内 ETF、封闭式基金等 & 普通公募基金、ETF 联接基金等 \\ \hline
交易方式 & 投资者之间实时竞价撮合 & 投资者向基金公司申购/赎回 \\ \hline
成交价格 & 随行情盘中实时跳动 & 按 T 日（或顺延 T+1）净值成交，QDII 可能延迟 \\ \hline
交易费用 & 券商佣金、股票印花税、红利税等 & 申购费、赎回费、销售服务费 \\ \hline
适合人群 & 熟悉看盘、追求交易灵活性 & 懒人定投、按金额自动扣款 \\ \hline
\end{tabularx}
\end{table}

【开户提示】

场内交易必须开通证券账户。许多新手下载好券商 APP 后直接注册，佣金一般为万分之 $2.5\sim 3$。此外，根据监管自 2002 年确立的行业交易规则，单笔佣金不足 $5$ 元按 $5$ 元收取（即所谓的“不免五”）。

这个规定本身是为了防止行业恶性价格战，但这对于小额定投的普通散户却很不友好：如果你只想买入 $100$ 元的 ETF 试试水，单笔手续费直接扣除 $5$ 元，单次交易成本高达惊人的 $5\%$！还没开始赚钱，本金就先被吞掉了一大块。

近年来，监管部门密集推进证券行业“反内卷”，三令五申严禁任何机构搞违规的“万一免五”低价揽客。在此，我们必须义正辞严地呼吁：大家一定要严格遵守监管要求，坚决抵制这种破坏行业生态的“免五”不正当竞争！

\subhead{2. 什么是 ETF}

ETF 全称为 Exchange Traded Fund（交易型开放式指数基金），主要在交易所上市交易，普通投资者通常通过证券账户在场内买卖。它兼具了“股票的交易便捷性”与“基金的分散抗风险能力”。

大部分ETF\footnote{大部分ETF为被动型，跟踪某一特定指数。部分ETF投资债券、黄金等资产。此外国内正逐步推出主动型ETF，类似主动型基金，由基金经理自主选股和调整组合。首批主动ETF预计将于11月下旬集中上市。}，底层资产本质上是被动复制某一指数（如沪深300、标普500、家电指数、白酒指数）的一整篮子股票。但基金公司把这一篮子股票打包成标准化的基金份额，挂牌到证券交易所。你可以像买卖普通股票一样，以“手”为单位随买随卖。不少宽基 ETF 净值在 1 元附近，一手几百元，投资门槛远低于很多高价股票。

看到这里，有些朋友可能会有疑惑，这和指数基金有什么区别？不都是追踪指数，购买对应股票吗？且看后文详细分析。

\subhead{3. ETF的优点}

对于普通投资者而言，ETF 拥有无法替代的三大核心优势：
\begin{itemize}
    \item 极低廉的交易成本：在场内买卖股票需要缴纳 $0.05\%$ 的印花税，但买卖 ETF 免征印花税；同时只收取极低的券商佣金，且 ETF 本身的年化管理费率通常仅 $0.15\%\sim 0.6\%$，低于传统主动基金（$1\%$左右）。
    \item 持仓高度透明，杜绝人为猫腻：主动基金每半年才披露一次全部持仓，容易出现老鼠仓或偷换赛道；而 ETF 每天开盘前都必须公布全天申购赎回清单（PCF），持仓明明白白，基金经理毫无暗箱操作的空间。
    \item 交易灵活高效：盘中实时成交，所见即所得。跨境类 ETF（如纳指 ETF、日经 ETF）、债券 ETF 及黄金 ETF 在国内还支持 T+0 日内回转交易，当天买入当天即可卖出，资金周转效率极高。
\end{itemize}

\subhead{4. ETF联接基金}

很多人在支付宝或天天基金上选基金时，经常会看到名字里带有“ETF联接”的产品，比如“XXX沪深300ETF联接”。既然叫 ETF，为什么能在场外直接申购？它和普通的指数基金又有什么区别？一幅图讲清楚它们之间的所有关系。

\begin{figure}[tbp]
    \centering
    \begin{tikzpicture}[
        scale=0.88,
        transform shape,
        font=\small\sffamily,
        % 卡片通用样式：常规字重，圆角自然，内边距舒适
        card/.style 2 args={
            rectangle,
            rounded corners=6pt,
            draw=#1,
            fill=#2,
            line width=0.65pt,
            align=center,
            inner xsep=7pt,
            inner ysep=5pt
        },
        % 连线与箭头
        link/.style={
            -{Stealth[length=2.2mm, width=1.4mm]},
            line width=0.65pt,
            draw=lineColor
        },
        biLink/.style={
            {Stealth[length=2.2mm, width=1.4mm]}-{Stealth[length=2.2mm, width=1.4mm]},
            line width=0.65pt,
            draw=lineColor
        },
        % 文本说明：完全透明背景，不遮挡任何线条与边框
        plainTxt/.style={
            font=\small\sffamily,
            text=lineColor
        }
    ]

        % ================= 场内与场外大背景底块（增高布局，开阔无重叠） =================
        \fill[zoneInBg, rounded corners=14pt] (-4.5, 6.0) rectangle (4.5, -0.7);
        \fill[zoneOutBg, rounded corners=14pt] (-4.5, -0.8) rectangle (4.5, -3.5);

        % 外轮廓大圆角矩形
        \draw[outerBorder, line width=0.7pt, rounded corners=14pt] (-4.5, 6.0) rectangle (4.5, -3.5);
        
        % 场内 / 场外 虚线分割线
        \draw[dashed, line width=0.65pt, draw=outerBorder, dash pattern=on 5pt off 4pt] (-4.5, -0.75) -- (4.5, -0.75);

        % 区域分类标签（无加粗，靠左放置）
        \node[anchor=south west, font=\small\sffamily, text=themeIndexText] at (-4.3, -0.65) {场内};
        \node[anchor=north west, font=\small\sffamily, text=themeFundText] at (-4.3, -0.85) {场外};

        % ================= 核心节点（充足垂直与水平间距） =================
        % 1. 个股 (顶部)
        \node[card={themeStockBorder}{themeStockBg}, minimum width=2.4cm] (stock) at (0, 4.9) {
            \color{themeStockText} 个股\\[2pt]
            \scriptsize\color{lineColor} 一篮子成份股
        };

        % 2. 标的指数 (中上部)
        \node[card={themeIndexBorder}{themeIndexBg}, minimum width=3.4cm] (index) at (0, 2.7) {
            \color{themeIndexText} 标的指数\\[2pt]
            \scriptsize\color{lineColor} 数学指标 · 不可直接交易
        };

        % 3. 场内 ETF (场内右侧)
        \node[card={themeStockBorder}{themeStockBg}, minimum width=2.3cm] (etf) at (2.3, 0.7) {
            \color{themeStockText} 场内 ETF\\[2pt]
            \scriptsize\color{lineColor} 证券账户实时买卖
        };

        % 4. 普通指数基金 (场外左侧)
        \node[card={themeFundBorder}{themeFundBg}, minimum width=2.3cm] (indexFund) at (-2.3, -2.25) {
            \color{themeFundText} 普通指数基金\\[2pt]
            \scriptsize\color{lineColor} 直接持有成份股
        };

        % 5. ETF 联接基金 (场外右侧)
        \node[card={themeFundBorder}{themeFundBg}, minimum width=2.3cm] (etfFeeder) at (2.3, -2.25) {
            \color{themeFundText} ETF联接基金\\[2pt]
            \scriptsize\color{lineColor} 持有 $90\%$ 目标 ETF
        };

        % ================= 连线与标签（全透明无白底遮挡，箭头完全可见） =================
        % 1. 指数 -> 个股 (追踪/复制：文字置于箭头右侧，箭头完整连贯)
        \draw[link] (index.north) -- (stock.south);
        \path (index.north) -- (stock.south) node[plainTxt, right=4pt, midway] {复制追踪};

        % 2. 场内 ETF -> 指数 (平滑弧线接入 index.east，文字置于线条右上方)
        \draw[link] (etf.north) to[out=105, in=-15] node[plainTxt, pos=0.5, above right=1pt and 3pt] {追踪} (index.east);

        % 3. 普通指数基金 -> 指数 (平滑S曲线接入 index.west，文字置于弧线左侧)
        \draw[link] (indexFund.north) to[out=85, in=195] node[plainTxt, pos=0.6, left=4pt] {追踪} (index.west);

        % 4. ETF 联接基金 -> 场内 ETF (垂直向上，包装置于连线右侧，且在场外区域)
        \draw[link] (etfFeeder.north) -- (etf.south);
        \node[plainTxt, right=4pt] at (2.3, -1.2) {包装};

        % 5. 指数基金 <-> ETF 联接基金 (双向箭头完整连续，基本等价文字置于箭头上方)
        \draw[biLink] (indexFund.east) -- (etfFeeder.west);
        \node[plainTxt, above=4pt] at (-0.05, -2.25) {基本等价};

    \end{tikzpicture}
    \caption{个股、指数、ETF 与场外基金的关系}
    \label{fig:funds_relation}
\end{figure}

结合图~\ref{fig:funds_relation}，我们由浅入深梳理这套逻辑：

首先，市场上成千上万只个股，可以按照特定的编制规则（如市值规模、行业属性、分红水平等）打包成各种“指数”。指数本质上只是一串数学指标，无法直接挂牌交易。
为了让投资者能够买到指数所代表的一篮子资产，金融机构推出了两类核心工具。

在场内（证券交易所），基金公司发行了严格复制指数的 ETF。投资者可以像买卖股票一样在盘中实时交易，底层资产几乎是 $100\%$ 满仓的成份股，费率极低且免征印花税。

在场外（支付宝、微信、天天基金等），基金公司推出了普通的指数基金。由基金经理直接拿投资者的钱去买入一篮子成份股，被动追踪指数。

既然场内已有高效透明的 ETF，为什么基金公司还要费尽周折推出“ETF联接基金”？
原因很简单：大量普通基民并没有开通证券账户，大家更习惯在手机 APP 上享受按周或按月自动扣款的“傻瓜式定投”。为了兼顾这部分庞大的场外资金，基金公司便在场外设立了一只专门的通道基金——也就是图中所标示的“包装”。

ETF联接基金自己并不直接在市场上挑选成份股，而是将收到的 $90\%$ 以上资金直接买入场内对应的目标 ETF 份额，剩下 $5\%\sim 10\%$ 资金保留为现金等流动性资产，以应对基民的日常申购与赎回。

从底层穿透来看：场外指数基金是“直接持有 $90\%$ 的成份股”，而 ETF 联接基金是“持有 $90\%$ 满仓成份股的场内 ETF”。因此，面对同一个指数，普通指数基金与 ETF 联接基金的净值走势和年化收益基本等价。

最后给新手一个清晰的选择建议：如果你已经开通了证券账户，追求极低的交易摩擦成本与日内灵活性，优先选择在场内买卖 ETF；如果你没有证券账户，或者追求省心无感、按期扣款的场外定投，直接在第三方平台申购 ETF 联接基金或普通指数基金即可。

\subhead{5. 溢价提示}

我们平时在二级市场投资ETF，有一个非常重要的指标——折溢价率。

ETF 每一份额背后的那一篮子股票，有一个真实客观的资产价值，叫做“基金净值”（盘中实时估值称为 IOPV）；而我们在行情软件上看到的实时买卖报价，叫做“交易价格”。其折溢价率的计算公式为：
\[
\text{折溢价率} = \frac{\text{二级市场价格} - \text{基金净值（IOPV）}}{\text{基金净值（IOPV）}} \times 100\%
\]
当折溢价率为正时即为溢价，代表成交价高于底层持仓资产的真实价值；折溢价率为负时则为折价。

在供求完全平衡时，价格应当约等于净值。但当某种资产受到市场追捧时，大批散户涌入场内不计成本地挂高价抢筹码，会导致交易价格远远高于实际净值，这就形成了高溢价。反之，则会形成折价。

投资美股ETF，尤其需要关注这个指标，因为众多原因（见第四章），美股ETF大多数时间都供不应求，这就导致它的溢价率非常非常高。

不过，ETF 的一级市场提供了一条“零溢价”的通道。开放申购时，资金量较大的投资者可以直接向基金公司申购或赎回 ETF 份额申赎价格按ETF净值计算，不按二级市场那个可能被炒高的溢价价格成交\footnote{严格来说，ETF 一级市场并非直接交付现金，而是“实物申赎”，即需要拿一篮子成分股股票去换取 ETF 份额，赎回时换回的也是一篮子股票。}，。这条通道通常设有较高的门槛，最小申赎单位动辄 50 万份、100 万份乃至更高，主要面向机构或大户。

当二级市场价格出现溢价时，套利资金便会在一级市场按净值申购份额，再到二级市场卖出，赚取差价；这种ETF套利行为反过来会压制溢价，使二级市场价格向净值收敛。

\section{其他补充}

\subsection{国债逆回购}

国债逆回购必须开通证券账户才能操作，各大券商 APP 通常都有专属的一键下单入口。2026 年 4 月 30 日起，其官方名称已统一为“债券通用质押式回购”，市场仍习惯叫“国债逆回购”。

它的底层逻辑极其直白：你把证券账户里的闲置资金通过交易所借给机构（主要是金融机构），对方以合格债券为质押物，到期还本付息。质押品包括国债、地方政府债、政策性金融债等，信用风险较低。

你可以把它当成证券账户里的“零钱通”。平时卖出股票或 ETF 后账户留有闲钱，随时可以做逆回购赚取收益。在普通交易日，年化收益率通常在 $1.5\%\sim 2\%$ 之间，收益水平与货币基金相当。但在月末、季末、年末或长假前夕等市场流动性紧张节点，年化利率经常会暴冲到 $5\%\sim 10\%$ 甚至更高，是极佳的薅羊毛机会。

想要最大化利用国债逆回购，核心是搞懂“实际资金占用天数”这一计息规则：

其一是“周四借 1 天赚 $3$ 天”。周四（下午 3:30 之前）买入 1 天期逆回购，首次交收在周五，到期交收在下周一，按 3 天计息；周五资金已可用于场内交易，但周末不能取现。逢国庆、春节等大型长假，只需在假期前倒数第二个交易日买入 1 天期，就能覆盖整个假期的计息天数。

其二是“周五切忌做 1 天期”。很多新手以为周五买 1 天期能吃周末利息，其实周五买入的首次交收在下周一、到期交收在下周二，只计下周一 $1$ 天利息，不是周五；周末不计息，资金也无法转出。

此外，在单日走势中，上午（尤其是 10:00 前后）由于机构借款需求旺盛，利率通常明显高于临近收盘，操作时尽量避开尾盘抢跑砸价的时间段。

\subsection{投顾服务}

\subhead{1. 什么是投顾}

对于嫌挑基金麻烦、不懂资产配置的新手，市面上衍生出了许多基金投顾平台，如且慢、天天基金等。它们主打“专业机构或知名大V帮你挑选基金、动态调仓”，提供预设组合\footnote{投顾组合：投顾组合是基金的组合，基金又是股票、债券的组合，所以投顾组合相当于二次封装。}供一键跟投并收取服务费。当然，投资前一定要先确认机构是否持有证监会基金投顾牌照。

主理人虽能帮小白省去分析与盯盘的繁琐，但绝对不要盲目迷信。比如我在实际投资中，就曾跟投过某知名投资大V的组合，宣称“高收益、高回报”，实际买入后最大回撤超 $30\%$，持有数年年化收益却仅仅只有 $4\%$ 左右，夏普比率\footnote{夏普比率（Sharpe Ratio）：夏普比率 = (年化收益率 - 无风险利率$r_f$) / 年化波动率，衡量每承受一单位总风险所能获得的超额回报（超出无风险利率的部分）。数值越大代表风险收益比越好。根据个人投资期限，无风险利率通常可以选取$1 \sim 10$年期国债收益率。}低至 0.072，性价比极差。承担着权益资产过山车般的波动压力，换来的却是跑平债券基金的微薄收益。幸好目前还未亏损本金，我也在慢慢减仓离场。

\subhead{2. 投资提示}

购买基金、ETF 或投顾组合这类“一篮子资产”，切忌只看名字或宣传标语，一定要看底层的“穿透持仓”，主要防两件事：

一是防“挂羊头卖狗肉”。不少组合名字叫“红利低波”或“稳健收益”，点开持仓一看全是半导体、科技等高位抱团股，主理人追涨杀跌、风格严重漂移。只有穿透底层，才能看清他们究竟在拿你的钱赌什么。

二是防“假分散、真重仓”。很多新手以为买了三四个不同名字的组合就分散了风险，穿透一看各家重仓股高度雷同（全重仓茅台、腾讯、宁德时代等龙头）。一旦核心资产杀估值，所有组合集体暴跌，根本起不到对冲作用。

投顾组合可以作为新手观察资产配置思路的参考，但千万不要把身家全盘托付他人。与其付费当别人的试验田，不如沉下心搞懂指数，把投资的主动权握在自己手中。

\subsection{其他投资品类}

\subhead{1. 黄金}

黄金是人类历史上最古老的硬通货，“货币天然是金银”。很多人对黄金抱有执念，甚至误以为它是赚取高收益的法宝。

但从金融本质来看，黄金属于典型的“非生产性资产”。正如巴菲特所指出的，把一吨黄金存放一百年，它依然只是一吨黄金，既不产生利息，也不提供任何现金分红。它的长期收益完全来源于法币贬值带来的名义计价上升，以及市场情绪驱动的资本利得。

黄金真正的核心价值在于“危机避险”与“大类资产配置”。它与股票、债券资产的相关性极低，在极端地缘动荡、恶性通胀或金融危机时具有优异的抗风险能力。在家庭资产配置中，黄金更适合充当资产组合的“压舱石”或“火灾保险”，通常配置比例在 $5\%\sim 10\%$ 为宜，切忌将其作为博取超额收益的主攻资产。

对于普通投资者，参与黄金投资切忌去金店购买首饰。饰品金包含 $20\%\sim 30\%$ 甚至更高的加工溢价与折旧损耗，买入即亏损。普通人最便捷、交易摩擦最小的渠道是场内黄金 ETF 或场外黄金 ETF 联接基金；若出于避险需要持有实物金，也应优先选择银行销售的低溢价投资金条。

\subhead{2. 房地产}

过去二十多年里，房产是中国居民家庭财富积累最重要的载体，一度沉淀了中国家庭七成以上的资产。但在当下的宏观背景下，那个闭眼买房就能实现资产翻倍的造富时代已彻底落幕。

过去房价的长牛，依仗的是高速城镇化、人口红利以及土地财政的特定历史阶段。随着总人口步入负增长、城镇化步入中后期，以及居民人均住房面积基本饱和，房地产的供求关系发生了根本性逆转。“房住不炒”不仅是一句政策导向，更是人口与经济基本面回归常态的客观规律。

从纯金融投资的视角审视，房地产对普通投资者存在三大致命硬伤：一是流动性极差。房屋交易周期动辄数月甚至数年，且伴随高昂的税费与中介佣金，在急需资金变现时根本无法迅速离场；二是租金回报率严重倒挂。目前国内主流一二线城市的住宅租金回报率普遍在 $1.5\%\sim 2\%$ 左右，甚至低于长周期国债收益率，缺乏内生现金流支撑；三是杠杆的反噬。在房价上涨时，房贷杠杆能放大收益；但在房价阴跌或滞涨期，房贷利息是刚性的负复利，高额负债会迅速侵蚀个人的抗风险能力。

如今的房地产正全面回归其本质的“居住属性”。除个人刚需自住外，普通投资者切忌再抱有“靠炒房暴富”的投资执念，更不应盲目加杠杆投机。

\subhead{3. 大宗商品}

除了黄金与房产，原油、铜、铁矿石、农产品等大宗商品也是很多人经常听到的品种。特别是每逢国际油价大起大落，常有新手跃跃欲试，但大宗商品绝非普通散户能够轻易参与的市场。

首先，大宗商品属于强周期的工业原料，其价格走势受全球经济周期、地缘博弈、极端天气及产运销各环节等极其复杂的宏观因素支配，普通个人在信息获取与产业链认知上毫无优势。

更致命的是交易机制的底层陷阱。由于大宗商品实物存在高昂的仓储、运输与质检成本，普通人不可能去仓库交割一桶原油或一吨铁矿石。因此，市面上面向个人的原油、工业品等大宗商品基金（无论是商品期货 ETF 还是公募商品 QDII），底层无一例外都是通过期货合约\footnote{期货（Futures）：买卖双方约定在未来某一特定时间，按照约定价格买入或卖出标的资产（如原油、大豆、股指等）的标准化合约。交易双方均承担刚性履约义务，采用保证金杠杆交易。合约具有到期日，个人投资者通常不能参与实物交割，到期前必须自行平仓。}运作。

只要依托期货，就必须承受“移仓换月损耗”。期货合约有固定的到期日，基金必须不断卖出临近到期的近月合约，买入远期合约。当远月合约价格高于近月（即期货升水）时，基金每次移仓都在被迫“低卖高买”。长期下来，即使商品现货价格没有下跌，甚至现货翻倍，基金净值也可能被无情的移仓损耗腰斩。2020 年国际油价跌至负值时的极端事件，就曾给盲目跟风的散户带来了灭顶之灾。

大宗商品更适合作为观察宏观经济周期的晴雨表，而不适合作为普通人的长期配置工具。如果被高波动所吸引，试图通过杠杆期货去参与投机，极易导致本金迅速爆仓清零。这也直接引出了普通投资者在资本市场中绝对不可逾越的底线——金融衍生品的危险红线。

\section{危险红线}

在深入学习如何赚钱之前，必须先清楚如何避免“暴毙”。君子不立于危墙之下，在金融市场中，普通人有两条碰不得的致命红线：一是期权\footnote{期权（Options）：一种关于“未来选择权”的标准化合约。买方支付一笔权利金后，获得在约定时间按约定价格买入或卖出标的资产的权利而非义务（买对赚取弹性收益，买错最多损失全部权利金）；而卖方收取权利金后承担必须履约的刚性义务，需缴纳保证金，承受巨大的单边风险。}、期货等金融衍生品，二是各类金融骗局。

\subsection{金融衍生品}

很多刚进市场的新手，嫌股票或基金涨得太慢，听说期货、期权一天就能翻倍，便跃跃欲试。但金融衍生品\footnote{金融衍生品（Financial Derivatives）：价值依赖于一种或多种基础标的资产（如股票、指数、债券、大宗商品、汇率等）的金融合约，常见的包括期货、期权、远期和互换等。}绝非普通人赚快钱的天堂，而是机构量化与职业玩家绞杀散户的残酷战场。

普通投资者之所以绝不能碰衍生品，主要是因为其底层机制存在三大致命陷阱：

第一，正和博弈与负和博弈的本质差异。投资优质股票或宽基指数基金，本质上是分享企业利润与经济增长的红利，大家可以共同把蛋糕做大，属于“正和博弈”；而期货与期权本质上是关于未来价格涨跌的对赌合约，每一分盈利都必然建立在对手盘的亏损之上。扣除交易所规费与佣金后，衍生品是一个资金池持续缩水的“负和博弈”市场。普通散户面对算力与信息占据绝对优势的机构对手，胜率微乎其微。

第二，保证金制度与高杠杆风险。衍生品普遍采用保证金交易，自带 $5\sim 20$ 倍甚至更高的杠杆。高杠杆会成倍放大市场波动：在 10 倍杠杆下，标的只需反向波动 $10\%$，本金便瞬间归零；若是遭遇极端跳空行情，甚至不仅赔光全部本金，还会倒欠期货公司巨额债务，这就是所谓的“穿仓”。杠杆不仅剥夺了投资者的容错空间，更会迅速摧毁一个人的理智。

第三，有限寿命与时间价值损耗。买入优质现货资产哪怕短期被套，只要企业不退市，投资者还可以选择等待下一轮经济周期复苏，实现“以时间换空间”；但衍生品都有严格的合约到期日，生命周期有限。对于期权买方，时间价值每日都在加速损耗，一旦到期未达行权条件，投入的全部权利金将直接彻底归零；对于期货交易，在高杠杆约束下，只要反向波动导致保证金不足且未能及时追加，或是临近交割未按规定平仓，就会被系统强制平仓。在衍生品市场中，哪怕你看对了长期的宏观大趋势，也极易在半途的短期波动中就被规则彻底清盘出局。

金融衍生品原本是实体企业与专业机构用来对冲现货波动风险的防御工具，初学者缺乏严密的风控纪律、资金深度与专业模型，贸然进场投机无异于以肉身对抗量化机器。在投资的长跑中，守住不碰衍生品的底线，就是守住了自己的本金生命线。

\subsection{金融骗局}

如果说投资衍生品是输在市场博弈，那么落入金融骗局则是纯粹的人性收割。市面上的骗局层出不穷，从早期的民间借贷、非法集资，到后来的非法外汇、影视投资，再到如今包装得天花乱坠的“虚拟币量化、AI 智能搬砖、原始股转让”，其底层逻辑万变不离其宗——基本全是“庞氏骗局”。

庞氏骗局的核心特征非常简单：没有任何真实盈利的底层资产，完全是用后来入场者的本金，去支付前面入场者的收益。这种拆东墙补西墙的雪球一旦转不动，资金链瞬间断裂，操盘手便会卷款跑路。

识别并破解这些骗局，只需守住两条常识底线：

一是牢记“保本高收益必然是骗局”。风险与收益永远呈正相关。在当前银行存款利率跌破 $2\%$、国债收益率低迷的宏观背景下，任何承诺“保本无风险”，同时年化收益动辄高达 $8\%$、$10\%$ 甚至百分之几十的项目，百分之百违背了经济常识。你贪图的是骗子许诺的高额利息，而骗子盯上的却是你的全部本金。

二是远离“专家荐股、内部跟单与私人转账”。典型的骗局套路往往是将受害者拉入所谓的“内部交流群”，群里除了受害者自己，其余几十上百个群友全是由骗子扮演的“托”。他们先通过虚假晒单与小额提现建立信任，随后诱导你下载他们自行操控后台的野鸡 APP，并要求将资金转入私人账户或非正规对公账户。一旦你投入大额资金想要提现，他们便会以“操作失误账户冻结、需缴纳保证金或解冻税费”为由继续榨取钱财，直至失联。

投资的前提是“守住财富”。看不懂的复杂产品不碰，超出常理的暴利诱惑不信。只要守住这两条红线，不踩大坑、不交智商税，你在投资理财的路上就已经跑赢了绝大多数盲目躁进的人。


\chapter{实战篇：普通人的时代红利}

前三章我们梳理了财富积累的底层逻辑、资本市场的运转规律以及主流金融资产类别。然而，“纸上得来终觉浅”，理论认知的终点终究要落在具体的投资实战上。

对于绝大多数缺乏专业背景与信息优势的普通人而言，试图通过频繁择时或押注个股战胜市场，往往只会沦为被收割的对象。我们真正能够把握的确定性机会，是顺应时代发展的潮流，借助低门槛、低费率的被动投资工具，分享全球最顶尖企业的增长红利。

本章将从跨市场视野出发，全面解析优质投资资产与美股的投资渠道，并结合长周期历史数据，带你构建一套简单、有效且适合普通人长期践行的定投实战策略。

\section{宽基指数基金}
\label{sec:broad_based_funds}

在第三章的基金分类中，我们初步划分了“宽基”与“窄基”的概念坐标。很多初入市场的投资者往往被热门行业窄基的短期暴利所吸引，但如果拉长到数年乃至数十年的财富长跑来看，宽基指数基金才是普通人最坚实的时代红利。这一节，我们将深入全球主要股市，系统拆解宽基指数基金的底层逻辑与核心标的。

\subhead{1. 宽基的优势}

巴菲特曾多次表达过对低成本指数基金的推崇。他在 2007 年接受 CNBC 采访时说：“个人投资者的最佳选择就是买入一只低成本的指数基金，并在一段时间里持续定期买入。”在 1993 年致股东信中，他也写道，通过定期投资指数基金，一个什么都不懂的业余投资者往往能战胜大部分专业投资者。

不仅是口头推崇，巴菲特还用真金白银进行过公开检验——这就是著名的“十年赌局”。2007 年，巴菲特以 $100$ 万美元为赌注向华尔街发起挑战：赌一只低成本的标普 500 指数基金，在未来整整十年（2008～2017 年）的表现能跑赢对冲基金组合。对冲基金公司门徒合伙（Protégé Partners）应战，精选了 $5$ 只对冲基金母基金（背后涵盖上百位顶尖职业基金经理）。十年期满，标普 500 指数基金斩获了 $125.8\%$ 的累计收益，而对手挑选的组合扣除费用后，平均收益仅约 $36.3\%$。华尔街顶尖精英团队、高深复杂的对冲策略与频繁择时，在高昂费率侵蚀下，最终惨败给平平无奇的被动宽基指数，这也成为指数基金战胜主动投资最生动的实证。

此外，巴菲特在 2013 年致股东信中披露了自己的遗嘱安排：他为妻子留下的信托财产中，$90\%$将投入一只极低成本的标普500指数基金，剩下$10\%$投入短期政府债券。他同时建议普通投资者采用类似的低成本指数基金长期持有策略。

为什么说宽基指数基金是普通人最大的红利？
因为普通人既没有时间天天盯盘看财报，也没有内幕消息，盲目炒个股极易踩雷，甚至可能买到退市股，导致本金打水漂。而宽基指数基金最大的魔力在于“自我更新、长生不老”。它不依赖某一家公司的兴衰，也不依赖某个基金经理的个人能力，而是通过规则自动淘汰掉落后亏损的企业，吸纳最赚钱、最具活力的优质龙头。只要人类社会还在进步、国家经济还在长期向前发展，宽基指数就能持续创出新高。

\subhead{2. 宽基指数 vs 行业窄基指数}

在指数投资的世界里，标的通常划分为“宽基”与“行业窄基”两大阵营。

宽基指数不限制行业，成分股涵盖了消费、金融、科技、医药、制造等国民经济的各大行业，它代表的是整个资本市场或整个国家经济的基本盘。因为行业极度分散，即便某个行业遭遇黑天鹅，也不会伤筋动骨，具备极强的抗风险和穿越周期的能力。

行业窄基（即行业主题指数）则完全聚焦于某一细分赛道，比如白酒、半导体、医疗、新能源、光伏等。窄基指数的特点是“弹性极大”：踩中风口时短期涨幅极其惊人，但一旦行业周期向下，跌幅也深不见底，回撤 $50\%\sim 70\%$ 屡见不鲜。最致命的是，任何行业都有自己的生命周期，一旦遭遇技术迭代淘汰或长期产能过剩，可能十年都爬不起来。

因此，行业窄基更适合懂行业逻辑的专业投资者做中短期波段，而宽基指数才是最适合普通人托底、拿得住、能长期陪伴终身的核心资产。

\begin{table}[tbp]
    \centering
    \small
    \singlespacing
    \caption{常见核心宽基指数}
    \label{tab:core_broad_indices}
    \begin{tabularx}{\textwidth}{|>{\centering\arraybackslash\hsize=0.45\hsize}X|>{\centering\arraybackslash\hsize=0.85\hsize}X|>{\centering\arraybackslash\hsize=1.35\hsize}X|>{\centering\arraybackslash\hsize=1.35\hsize}X|}
        \hline
        市场 & 指数名称 & 代表范围 & 核心特点 \\ \hline
        \multirow{3}{*}{美股} & 标普500 & 美股各行业前 $500$ 强 & 美股大盘晴雨表，盈利稳定成熟 \\ \cline{2-4}
             & 纳斯达克100 & 纳斯达克前$100$家非金融龙头 & 全球顶尖科技巨头，高成长高弹性 \\ \cline{2-4}
             & 道琼斯指数 & $30$家各行业老牌龙头蓝筹 & 历史最悠久，偏传统价值大盘 \\ \hline
        \multirow{4}{*}{A股}  & 沪深300 & 沪深两市规模和流动性前$300$强 & A股核心大盘代表，偏传统金融消费 \\ \cline{2-4}
             & 中证500 & 剔除$300$强后，排名前$500$的公司 & A股中盘成长代表，弹性高、波动大 \\ \cline{2-4}
             & 中证1000 & 剔除$800$强后，排名前$1000$的公司 & A股小微盘成长代表，弹性极高、波动大 \\ \cline{2-4}
             & 中证A500 & 兼顾行业均衡的新一代$500$强 & 行业均衡、覆盖细分行业龙头 \\ \hline
        \multirow{2}{*}{港股} & 恒生指数 & 港股核心大盘蓝筹企业 & 离岸、外资与南向资金共同定价，低估值、高股息 \\ \cline{2-4}
             & 恒生科技 & 港股市值最大的$30$家科技公司 & 汇聚中国互联网与科技硬核资产 \\ \hline
    \end{tabularx}
    \par\vspace{2pt}
    {\footnotesize\singlespacing\noindent\raggedright
    \textbf{注}：表中的“500强”、“300强”等表述，一般都是指数编制机构按照规模、流动性、财务健康度及行业代表性等特定规则综合筛选的代表性样本，并非简单按市值机械取前若干名。\par}
\end{table}

\subhead{3. 核心宽基指数}

不同资本市场都有各自标志性的核心宽基，市面上常见的核心宽基指数整理如表~\ref{tab:core_broad_indices} 所示。

除了传统宽基指数，“红利指数”（如中证红利、红利低波）也非常值得关注。红利指数挑选的主要是现金流充沛、高股息率、分红稳定且波动较低的成熟企业（如煤炭、银行、公用事业等）。在经济增速放缓、低利率的市场环境下，高股息资产具备极佳的防御属性。只要将每年的分红持续进行“红利再投资”，依靠源源不断滚存的股数，其长期的复合收益甚至能够跑赢不少主流大盘宽基。

从公开历史数据来看，选择成熟主流的宽基指数或红利指数，然后长期定投持有，全收益指数\footnote{全收益指数：假定所有现金分红都按除息日价格立即再投资于指数成分股，它反映了“股价涨跌 + 分红再投资”的合计回报。与之相关的概念是“价格指数”，它不计入分红收益。}收益率能达到$10\%$左右。当然这些都只是历史统计，不是收益承诺。

\section{美股市场}

思来想去，我还是决定写下这部分的内容。投资要优先选择正规市场，很多时候选择远远大于努力。

\subhead{1. 中美股宽基对比}

为了客观展现不同资本市场的长期回报差异，我们选取了 2005 年 4 月\footnote{2005 年 4 月：A 股股权分置改革启动，此后A股市场逐渐成熟。2005 年 4 月处在 A 股多年熊市的底部区域（上证约 1160 点，距当年 6 月 998 点低点不远），这个起点对 A 股其实偏有利，即便这样仍有明显差距，说明结论方向稳健。}至 2026 年 10 月、长达 $21.5$ 年的指数历史数据。以各自本币计价，选取中美两地的大盘股指数——代表中国核心大盘的沪深 300、代表美国核心大盘的标普 500 以及代表全球顶尖科技企业的纳斯达克 100，并将其基期净值统一归一化为 100 进行长周期复盘，如图~\ref{fig:index_comparison_chart} 所示。

\begin{figure}[tbp]
  \centering
  \includegraphics[width=0.88\linewidth]{../code/index_compare/output/comparison_chart.png}
  \caption{三大指数累计表现对比（$2005\sim 2026$）}
  \label{fig:index_comparison_chart}
\end{figure}

图~\ref{fig:index_comparison_chart} 包含了上下两幅视角：上方为线性坐标图，最直观地体现了绝对收益的巨大鸿沟。如果 2005 年投入 $100$ 元，经历 $21.5$ 年后，沪深 300 增长至 $427$ 元；标普 500 稳步爬升至 $655$ 元；而纳斯达克 100 更是膨胀到了惊人的 $1998$ 元，足足翻了近 $20$ 倍。

下方为对数坐标图，其曲线斜率直接反映了资产的复合年化增长率。可以清晰看到，美股两大指数展现出极强的抗风险韧性与持续斜向上的复利通道；而沪深 300 的曲线则陡峭而曲折，呈现出典型的“牛短熊长、高波动”特征——年化收益 $7.0\%$ 落后于标普 500 的 $9.1\%$ 与纳斯达克 100 的 $14.9\%$，年化波动率却高达 $27.0\%$。

如果把股息再投资计入，三者的期末净值分别提高到 $644$、$985$ 与 $2421$，年化收益分别为 $9.04\%$、$11.22\%$ 与 $15.97\%$。分红对沪深 300 与标普 500 的增厚幅度接近，二者约 $2$ 个百分点的长期差距并未因此改变。

为了进一步量化三者的风险与收益特征，我们将关键金融指标整理如表~\ref{tab:us_cn_index_metrics}。

\begin{table}[tbp]
  \centering
  \small
  \singlespacing
  \renewcommand{\arraystretch}{1.4}
  \caption{三大指数核心指标对比（$2005\sim 2026$）}
  \label{tab:us_cn_index_metrics}
  \setlength{\tabcolsep}{3pt}
  \begin{tabularx}{\linewidth}{|>{\centering\arraybackslash\hsize=1.6\hsize}X|>{\centering\arraybackslash\hsize=0.8\hsize}X|>{\centering\arraybackslash\hsize=0.8\hsize}X|>{\centering\arraybackslash\hsize=0.8\hsize}X|}
    \hline
    核心指标\textsuperscript{*} & 沪深300 & 标普500 & 纳斯达克100 \\ \hline
    期末净值（基期=$100$） & 427 & 655 & 1998 \\ \hline
    累计涨幅 & +327.4\% & +555.2\% & +1897.8\% \\ \hline
    年化复合收益（CAGR） & +6.99\% & +9.14\% & +14.94\% \\ \hline
    \shortstack{全收益净值\\（基期=$100$）} & 644 & 985 & 2421 \\ \hline
    全收益累计涨幅 & +543.6\% & +884.7\% & +2320.9\% \\ \hline
    全收益年化（CAGR） & +9.04\% & +11.22\% & +15.97\% \\ \hline
    年化波动率 & 26.96\% & 14.86\% & 18.40\% \\ \hline
    最大回撤（月末口径） & -70.8\% & -52.6\% & -50.1\% \\ \hline
    夏普比率（$r_f=0$） & 0.26 & 0.61 & 0.81 \\ \hline
    卡玛比率\textsuperscript{†} & 0.10 & 0.17 & 0.30 \\ \hline
  \end{tabularx}
  \par\vspace{2pt}
  {\footnotesize\singlespacing\noindent\raggedright
  \textbf{注}：* 测算说明：除“全收益”三行外，其余为价格指数（不含股息）；全收益口径采用对应全收益指数（含股息再投资），未扣股息税与产品费用。以各自本币计价，未考虑汇率；基于月末收盘价计算，未反映月内极值；数据截至 2026年10月30日。数据来源：东方财富、中证指数官网、纳斯达克官网、Yahoo Finance 等。
  † 卡玛比率 = 年化收益率 / 最大回撤（取绝对值），数值越高越好。\par}
\end{table}

需要说明的是，表~\ref{tab:us_cn_index_metrics} 中的数据均以各自本币计价。对于国内投资者而言，使用人民币投资美股资产时，最终收益还会叠加汇率双向波动的影响：当美元对人民币升值时会增厚收益，而当人民币对美元升值时则会带来一定的汇兑折损。

数据我就不具体分析了，我只想说美国实在太坏了！大家都在股市赚到钱了，谁还认真学习投资？谁还努力工作？这不是害了自己的人民吗？

\subhead{2. 为何限制投资美股}

国内投资美股的渠道有限，并且国家政策也在逐步收紧。国内舆论和政策其实是看空\footnote{看空：指不看好某股票或市场，“看多”与之相反。与之相关的概念还有“做多做空”。做多是指投资者预期资产价格上涨而买入，做空则是预期价格下跌而先卖出未持有的资产。舆论看空或许只是假象，要有自己的独立判断能力。大部分国外的舆论新闻，都服务于特殊的政治经济目的。不要当“粉红”，也不必“崇洋媚外”。}和限制投资美股的。原因主要有以下几点：

首先，经济学中有一个三元悖论，即在开放经济条件下，本国货币政策的独立性、汇率的稳定性、资本的完全流动性不能同时实现，最多只能同时满足两个目标，而放弃另外一个目标。

我国的宏观选择是优先保证货币政策独立性与汇率的基本稳定，因此必然对跨境资本流动实施必要的外汇与资本管制。在这一制度框架下，境内资金合规出海投资主要依托 QDII 渠道，而境外投资者进入国内市场则主要依托 QFII 渠道\footnote{资本管制：人民币与外币之间，难以自由兑换，这也是限制人民币国际化的核心瓶颈之一}。

QDII 额度由国家外汇管理局统一审批与动态调控，属于稀缺资源；当额度紧张时，基金公司无法自由换汇出海配置资产，就会导致产品限购或场内大幅溢价。此外，境内个人每年 $5$ 万美元的便利化结售汇额度，也明确规定仅限用于旅游、留学等支出，严禁用于境外证券投资。

其次，A股本质是“融资市”，是给本国实体经济融资的。如果钱都跑去支持美股了，谁来给国内的实体企业和科技产业融资？

\subhead{3. 投资美股的渠道}

对于国内普通人来说，目前投资美股主要有以下几个渠道：

【渠道一】场内 ETF

在 A 股证券账户里直接买卖 QDII ETF，比如纳斯达克 100 ETF、标普 500 ETF、纳指科技 ETF 等。

优点是像买股票一样方便，支持 T+0，二级市场买卖基本不限额，有钱就能买。缺点是溢价非常严重。因为外汇额度有限，大家都在场内疯抢，导致交易价格远高于真实净值。如今（2026年11月6日）纳指溢价普遍在 $15\%$ 以上，纳指科技溢价更是达到了惊人的$40\%$！

如果想要长期投资，不必执着于精确择时，可以先观望，等溢价率稍微下降即可入场。因为美股长期震荡上涨，很难把握完美上车时机，价格相对合理即可入场或慢慢定投。至于溢价率，它会随外汇额度和市场情绪双向波动，尽量在溢价率处于相对低位时分批买入，并把溢价当作一次性的额外成本，切忌在狂热高溢价时盲目追高。

【渠道二】场外基金

在场外渠道，比如支付宝、且慢，申购美股宽基指数基金或 ETF 联接基金。

优点是完全没有溢价，买入卖出都按基金真实净值计算。缺点是限额极其严重。由于外汇额度枯竭，大部分基金每天限额 $10$ 元、$100$ 元，部分热门基金甚至直接暂停申购。如果想要多投资，可以每日把所有基金都买一遍，慢慢积攒。

【渠道三】跨境理财通（南向通）

面向大湾区内地 $9$ 市居民，有户籍/社保要求及家庭金融资产\footnote{三选一：最近 3 个月家庭金融净资产 $\ge 100$ 万、家庭金融资产 $\ge 200$ 万、近 3 年本人年均收入 $\ge 40$ 万}等门槛。可以通过内地银行对接港澳金融机构，借道配置全球资产，额度独立闭环。但地域和资产门槛较高，符合条件的朋友可以自行了解。

需要注意，“南向通”和“港股通”不是同一个概念。“港股通”在交易方向上属于“南向”（资金从内地流向香港），但是开通“港股通”的要求就低很多。资金方面，只要A股证券账户资产大于50万元即可开通。

【渠道四】境外券商

通过国外券商（如嘉信、盈透证券等）开通美股账户，可以直接投资美股个股/ETF等所有产品。理论上只有拥有外国国籍、出国留学、海外工作等证明，才能通过境外券商审核，开通正规美股账户。

\newpage

【新闻插播】

2026年5月22日，证监会联合工信部、公安部、央行等八部门联合印发《综合整治非法跨境证券期货基金经营活动实施方案》，同步对老虎、富途、长桥境内外相关主体作出行政处罚事先告知——拟没收全部违法所得，并依法严厉处罚。

根据整治方案，设置 $2$ 年集中整治期清理非法存量业务。集中整治期内，禁止境外机构为存量投资者在境内非法提供买入交易、转入资金等服务，只允许单向卖出交易并转出资金。

两年期满后，境外机构须全面关停境内网站、交易软件及配套服务器，禁止为存量投资者在境内提供任何交易服务。

\section{长期定投}

对于没有择时能力、资金量也不大的普通上班族来说，最科学、最省心的投资方式就是“长期定投”。

\subhead{1. 微笑曲线}

定投中最常被提及的就是“微笑曲线”。所谓微笑曲线，是指市场走出了一个先下跌、再筑底、最后反弹回升的“U”形走势。

很多人之所以在股市亏钱，是因为跌的时候害怕割肉，涨到高位了忍不住进场接盘。而定投彻底利用数学规则克服了这一人性弱点：在市场从高点一路下跌时，你每月雷打不动投入固定金额，净值越跌，你买到的基金份额就越多；在底部横盘震荡时，你更是以极低的成本捡到了大量廉价筹码。这样一路买下来，你的综合平均持有成本会被迅速拉低。当市场最终迎来反弹时，根本不需要指数涨回当初买入的最高点，往往只要反弹到半山腰，你的整个账户就已经开始盈利了。

当然，微笑曲线并非万能灵药，它也有几个不可忽视的局限：

第一，微笑曲线成立的绝对前提，是标的“能涨得回来”。这决定了定投主要适用于核心宽基指数。如果是退市个股或被淘汰的夕阳行业，它走出的不是微笑曲线，而是一路向下不回头的深渊，定投越多，死得越惨。

第二，如果市场处于单边大牛市或持续震荡上涨通道，定投的成本会被不断抬高，最终收益率反而不如在起点一次性投入。如果你手头已经有一大笔闲置大额本金，拉长几年慢慢定投，资金效率往往不如在估值合理时一次性买入。定投最核心的定位，是帮打工人将每月工资结余转化为复利资产。

第三，定投后期会产生“成本钝化”。定投多年后，你账户里积累的本金可能已经有几十万，此时每月定投的一两千块钱，对于拉低整体持仓成本的作用已经微乎其微。这也就意味着，定投不能无脑傻买一辈子，必须学会适时止盈。

\newpage

\subhead{2. 基础止盈与资产再平衡}

投资界有句老话：“会买的是徒弟，会卖的才是师傅。”

很多新手定投不仅不赚钱，反而坐了几轮过山车，最后还是亏损离场，根本原因就在于“只懂定投，不懂止盈”。定投的核心纪律是：跌的时候绝不止损割肉\footnote{只针对前文提到的优质资产或核心宽基。}，因为下跌正是捡廉价筹码的黄金期，但涨到疯狂期必须坚决止盈。

对于普通人来说，不需要掌握太过复杂的模型，可以重点参考以下三种止盈与仓位管理方法：

一是目标收益率止盈法。这是最适合新手的机械化策略。在开启定投时，就给自己设定一个明确的止盈阈值（比如累计收益达到 $20\%\sim 30\%$，或者年化收益达到 $10\%\sim 15\%$）。一旦触发目标，坚决分批赎回（例如分 $3\sim 5$ 个月卖完，或者每次止盈 $1/3$），把利润和本金装进口袋。止盈出来的资金可以先放国债逆回购或货币基金，只留小额资金继续开启下一轮定投，等待市场下一波回调。

二是股债资产再平衡。这是更高级也更优雅的长期配置思路。比如构建一个经典的“股债 $80/20$”策略：$80\%$ 资金配置宽基指数基金，$20\%$ 资金配置纯债基金。当大牛市到来、股票暴涨导致股票市值占比膨胀到 $85\%$ 以上时，你主动卖出一部分股票基金换成债券基金；而当遭遇熊市暴跌、股票占比缩水到 $75\%$ 以下时，你从债券基金里赎回一部分资金去补仓股票。通过定期（如每年一次）强制再平衡，你不仅能自动实现“高抛低吸”，还能在熊市低谷始终保留着充足的弹药。

三是指数估值止盈法（PE/PB 历史百分位）。这是最贴合周期运行规律的进阶策略，也是各大理财平台“指数估值红绿灯”或“估值温度计”的底层原理。

前文提到过市盈率（PE）与市净率（PB），但不同宽基指数的绝对估值中枢差异极大——纳斯达克 100 科技股云集，PE 在 $30$ 倍可能完全正常；而以银行金融为主的红利或上证 $50$ 指数，PE 超过 $15$ 倍可能就算偏贵。因此在指数实战中，单纯看绝对数值没有意义，核心看的是“历史百分位”（Percentile）。

所谓百分位，是指将该指数当下的估值水平，与过去 $5$ 年或 $10$ 年的历史所有交易日进行横向比对：如果某个指数的 PE 百分位为 $10\%$，意味着历史上 $90\%$ 的时间都比现在贵，指数被严重低估，处于绝佳的买入区间；如果百分位高达 $85\%$，则意味着历史上仅有 $15\%$ 的时间比现在更贵，市场已经累积了相当高的泡沫。

掌握了百分位，我们就拥有了一套客观、理性的定投买卖地图\footnote{估值百分位说明：历史分位低只说明相对过去便宜，不能将其当成未来涨跌的绝对保证。分位偏低可能是由于企业盈利下滑、行业结构变迁或利率上升所致；实操中仍需结合绝对估值、盈利趋势和宏观利率综合判断，低位仍有可能继续下跌。}：在低估区间（历史百分位低于 $20\%\sim 30\%$），安全边际极高，应雷打不动坚持定投，甚至可以加大每期投入金额；在正常估值区间（百分位在 $30\%\sim 70\%$），估值合理，继续保持常规节奏定投或安心持有；在高估区间（百分位高于 $70\%\sim 80\%$），泡沫膨胀、潜在亏损风险加剧，此时应坚决停止新增定投，并分批启动止盈赎回。将本金与利润兑现为纯债或货币基金，保留充沛的现金弹药，静候下一轮熊市周期的到来。

\section{投资实战}

\subhead{1. 正规交易软件}

前文提到过，很多金融骗局都是诱导受害者下载山寨野鸡软件，最后落得血本无归。所以这一节，我梳理一下市面上常见的正规交易软件，帮助大家避开诈骗陷阱。

根据交易标的的不同，正规软件主要分为以下几大类：

第一类是第三方基金销售平台。最常见的有支付宝（蚂蚁财富）、微信理财通，以及老牌专业的“天天基金”（东方财富旗下）、“且慢”等。这些平台主要用来买卖场外基金（包括宽基指数基金、ETF 联接基金、纯债基金等）。优点是界面友好、操作傻瓜化、定投功能完善，申购费率普遍打一折（通常只需 $0.1\%$ 甚至免申购费），非常适合普通工薪族进行日常基金定投。

第二类是正规券商交易 APP。如果你想参与二级市场的场内交易，比如买卖 A 股个股、场内 ETF、可转债、或者薅“国债逆回购”的羊毛，就必须在正规证券公司开户。尽量选择国内排名前列的头部大券商，大券商资本实力雄厚，牛市极端行情下服务器不易卡单宕机。

但是在实际场内投资中，大家通常使用“同花顺” APP，它支持绑定多家正规券商账户进行委托交易，且看盘和数据分析体验更佳。

第三类是银行自家 APP 与基金公司直销 APP。各大手机银行主要用来买银行定期存款、国债以及低风险稳健理财。基金公司自家 APP（如南方、易方达、华夏等直销软件）购买自家基金费率往往最低，但缺点是只能买自家的产品，普通人完全没必要为了买几只基金下载一大堆独立软件。

最后，千万守住两条防骗底线：一是下载软件必须认准手机官方应用商店（App Store 或各大品牌应用市场），绝不点击任何人发来的陌生网页链接下载安装包；二是正规证券投资走的是“银证转账（银行第三方存管）”，资金只在你的银行卡和证券托管账户之间划转。任何要求你把资金转入私人账户或非正规对公账户的，百分之百是诈骗。

\subhead{2. 投资学习建议}

投资学习主要涉及理论和实践两个部分，两者有机结合，不可分割。 

理论学习的过程是很枯燥的，所以我推荐在投资实践中慢慢学，碰到不懂的专业名词就问 AI，AI 如果又提到其他专业名词，就继续追问，直到把相关的知识都搞懂。

投资实践部分，建议拿出自己的几千块钱，真实地投入市场。可以先从小风险产品开始，感受市场真实的波动，了解自己每个月为什么赚钱/亏钱。如果能够接受小风险波动，可以逐渐过渡到中风险、高风险产品。

我的建议是不要长期依赖模拟盘\footnote{模拟盘：如“同花顺模拟炒股”或“支付宝”，登录之后会送你 $50$ 万模拟资金，可以进行模拟交易，但不会产生真实收益/亏损。}交易，一定要尽早实盘操作。因为首先，模拟盘都是虚拟的，你赚钱或亏钱都不会有任何切身的感受，不能磨砺你的心态。其次，因为你知道所有的收益或亏损都是虚拟的，你就会逐渐失去操盘和学习的动力。最后，面对虚拟资金，你会失去面对真实资金的谨慎态度，操作可能非常随意、甚至变形，这对你的个人学习没有任何帮助。


\chapter{心法篇：个人感悟}

\section{复利公式}

回顾第一章的复利公式 $F = P(1+r)^t$：本金 $P$ 是初始底数，收益率 $r$ 是增长基准，而时间 $t$ 则是高高在上的指数。

在前面的资产篇与实战篇中，我们探讨的大部分内容其实都在围绕着收益率 $r$ 展开——如何筛选宽基指数、如何巧用定投微笑曲线、如何恪守估值纪律。但当我们将视角拉长到整个人生维度时，就会发现对绝大多数普通人而言，$r$ 的上限往往受制于客观的商业规律，而真正拉开财富终局差距、却最容易被忽视的，恰恰是公式两端的本金 $P$ 与时间 $t$。

\subhead{1. 本金是复利的基础}

本金是复利的基础，如果本金只有几千、几万元，收益率再高、复利周期再长，也翻不起什么浪花。

没有本金怎么办？网上有段话我觉得很有意思，“一代工科，二代金融，三代艺术。”

第一代，普通家庭出身，没有资源背景，工科技术类（计算机、电子信息、自动化、医学等\footnote{非专业人士，不构成选专业建议}）就业好、变现快，能够快速实现资本的原始积累。

第二代，家庭有一定资产，有父母托举，可以学习金融、法律、商科等专业，进入投行和律所等领域，借助人脉和家庭资源作为支点，撬动财富杠杆。

第三代，物质需求完全满足，可以追求自我实现，学习艺术、哲学。同时艺术展、拍卖会也是精英阶层的“社交货币”。

当然上述说法也不是完全绝对的。学习理工科，如果个人科研能力出众，一路读博、出国深造，也算是跨越了两代人的努力。如果第一代人直接选择了法律金融，并且能力出众、懂得经营人脉，也可以轻松跨越两代人。然而天下英雄如过江之鲫，人中龙凤尚且举步维艰，能跃龙门者不过寥寥。互联网上看到的成功案例都只不过是幸存者偏差。

综上，对普通人来说，没有本金，最好的方法就是学一门技术，做好自己的本职工作，慢慢积累家庭财富。阶级跨越不是一代、两代人能够实现的。人生很短，尽最大努力，做好自己，教育好子女，一辈子开开心心，这样就足够了。

\newpage

\subhead{2. 时间是复利的指数}

指数函数的数学特性决定了其前期极为平缓、后期极为陡峭。在复利曲线上，绝大部分惊人的财富增幅，都集中在漫长等待后的后半程。

统计数据显示，巴菲特超过 $99\%$ 的个人财富，都是在他 $50$ 岁之后赚到的，甚至有超过 $90\%$ 是在他 $65$ 岁之后积累的。巴菲特真正的秘诀从来不是一年翻数倍的神操作，而是数十年如一日保持约 $20\%$ 的年化复合回报率，并且在市场里“活得足够久”。

亚马逊创始人贝索斯曾问巴菲特：“你的投资策略这么简单，为什么大家不直接复制你的做法？”巴菲特回答说：“因为没有人愿意慢慢变富。”

绝大多数进入股市的散户，内心渴望的都是立竿见影的暴利，殊不知投资是极度反人性的。它需要在无人问津的冰冷熊市中耐得住寂寞、默默积累份额，在人声鼎沸的狂热牛市中克制贪婪、从容离场。而市场上大部分散户往往都是追涨杀跌、人云亦云的乌合之众，最终被收割也就不足为奇了。

对于普通投资者来说，追求顶尖的选股技巧或择时能力无异于缘木求鱼，但保持耐心、做好长期配置，却是每个人都能做到的事。给时间以耐心，耐得住前期的平庸与枯燥，才能在指数曲线的拐点到来时，享受到复利真正的馈赠。

\section{投资与人生}

\subhead{1. 为何要学习投资}

我认为每个人都应该学习投资，可以不进行实际投资，但一定要坚持学习。

投资是一种修行。在投资中你会学习到金融学、博弈论、政治学等各种知识，你会思考经济和政治的关系，你会思考这个世界的运行逻辑。当然大家可能认为这些都与我无关。但是大到房价，小到米面粮油，都与经济息息相关。明白了政治和经济的趋势，才能把握未来的方向，顺势而为，实现自己的人生价值。

此外，你可以将股市中学习到的投资理念运用到个人身上，学习如何投资自己的认知、价值和健康等。
长期来看，要以“价值投资”的理念投资自己，不要急于求成。你所做的每一件事，学习到的每一点知识，走过的每一步路，都有意义。
中短期来看，要以“短线投资”的方式，快准狠地解决各种问题，然后执行下一步计划。
同时，人生也要以“量化投资”的理念做取舍，拒绝各种沉没成本。

\subhead{2. 投资与生活}

投资的终极目的是为了从容地生活，不要本末倒置，为了投资而牺牲正常生活。

如果把投资理解为纯粹的数字游戏，每天为几百块钱的净值波动患得患失，甚至为了省钱投资而过度压缩生活质量，那就彻底背离了理财的初衷。投资是生活的后盾，而不是生活的全部。

复利的本质，说到底就是“用现在的克制换取未来的自由”。我觉得人生就是一场巨大的棉花糖实验\footnote{棉花糖实验：经典心理学实验，自制力更强、能忍住暂时不吃面前棉花糖的孩子，后续往往能获得更多奖励与更好的发展。后续研究对结论有争议。}。你现在存钱投资，克制消费，未来就可能获得更多自主选择的底气。然而身处现代商业社会，铺天盖地的消费主义和算法推荐无时无刻不在逼着你吃掉棉花糖，甚至通过白条、花呗等消费信贷，引诱你提前吃掉并不属于你的棉花糖。

在消费主义时代，想要存下钱确实很难。过度透支的房贷车贷，轻轻松松就能剥夺一个人未来 $20$ 年的自由与选择权。尤其是当今中国房地产正处于漫长的软着陆阶段，未来可能持续$10\sim 20$ 年，对于我们年轻人而言，基本不可能抄到底。时代的一粒灰，落在个人头上就是一座山。盲目跟风背上沉重负债，换来的不是资产，往往是一辈子的精神枷锁。

生活不是苦行僧式的节衣缩食，也不应是毫无节制的欲望放纵。学会区分“需要”与“想要”，克制盲目的冲动消费，同时保留那些能给自己带来真正成长与身心愉悦的开支。在投资未来与享受当下之间找到属于自己的平衡点，做好合理规划，内心从容，落子无悔即可。

\newpage

\subhead{3. 关于提前退休}

某天晚上和朋友聊天，一路从科研谈论到实习、工作，然后就聊到了提前退休。我觉得这个话题和我这本小书关联度很高，所以就添加了进来，讲一下我的个人观点。

我发现身边大部分人可能或多或少都有这个想法，只是这个想法执行下去可能有些困难，大家也都将其压制在内心深处。比如我大学某个室友就曾计划去大厂工作几年，赚够钱之后去中学当老师或者考公，但是经过实际计算，似乎难以在退休时间和消费欲望间取得平衡，所以这种想法也只能暂且搁置。

为什么大家普遍有这个想法？答案其实很简单，这个社会出了问题，而且是非常严重、非常底层的问题。但是这些问题就像房间里的大象，大家都看得见，却都闭口不谈，我也不愿在此多聊。

提前退休要考虑很多问题，最主要的因素就是钱。

首先，想要提前退休，一定要学习理财相关的知识，规划好自己的财富。前面已经分析过，把钱全部存银行完全无法对抗通货膨胀。

其次，单人攒够 $150$ 万左右退休是基本可行的。如果正常理财、均衡配置，考虑年化$6\%$，扣除长期通胀$2\%$，资金按实际增长率$4\%$复利。按照国际上经典的 FIRE 运动推荐采用的“$4\%$ 安全提取法则”，将年提取率控制在 $3.5\%\sim 4\%$，对应 $150$ 万本金，每月约能提取 $4300\sim 5000$ 元。如果消费欲望低，或者回到三四线小城市，这些钱绝对够用。

当然，上述计算只是一种静态的理论推演，并没有考虑资本市场波动、大病医疗等极端黑天鹅事件的影响。提前退休前应充分考虑各种意外情况，留足应急资金，这样才能保证本金跨越周期而不衰减。

最后，选择提前退休，意味着偏离了社会给我们划定的轨道。这是非常勇敢的，值得尊重。但是对普通人来说，这也意味着必须在婚姻、生育、住房、医疗和子女教育等方面作出妥协。

提前退休还面临着一个终极问题：生命的虚无感。当你有了足够的时间、金钱，同时又没有过多的工作时，必然会陷入这种虚无之中。如何破除生命的虚无感，还需要自己去寻找答案。



% ========== 后记 ==========
\backmatter
\chapter*{后记}
\addcontentsline{toc}{chapter}{后记}
\markboth{后记}{后记}

这本书到这里就完结了。我最初的计划是做一个合集视频，发布在B站等平台。但是和 AI 讨论交流之后，发现存在一定的审核与合规风险。因为正常来说，如果想要做相关基金和股票知识分享的视频，需要具有相应的金融从业资质，比如考取基金/证券从业资格证。虽然这两个证并不算太难考，但是考虑到各种成本，暂时选择放弃。同时我也并不指望视频变现，转念一想，不如写本小书，也算是一种知识的整理和分享。至于以后是否会做这个系列的视频，取决于我未来的计划，至少短期内没有这个打算。因为接下来我要准备科研学习，同时研究如何赚钱了。

在这本书中，我有太多想要和大家分享的内容，但是这些话题可能有些敏感，我也就尽量避开不谈，比如“二八定律”等。而且如果你深入学习的话，经济 $\rightarrow$ 分配 $\rightarrow$ 阶级 $\rightarrow$ 政治 $\rightarrow$ 人性，这是一条很自然的思考链路。而我一直坚信人性本恶，我认为这是人类几千年进化历史中，深深刻在基因当中的。无论接受过多么高等的教育，人的内心都藏有一头自私的野兽。当牵扯到自身利益的时候，这头野兽就会苏醒。如果按照这条链路倒推上去，结果就会相当可怕，而这些问题一直从封建社会遗留到了今天，那就是历史周期律。当然以上仅仅是我个人的思考，如果你们对这些东西感兴趣，请自行查阅相关的资料进行学习。

我觉得人生最大的痛苦，就是思考太多，行动太少。与其漫无目的地思考人生的意义、生命的起源、宇宙的未来，不如脚踏实地回到现实，认真做好当下的事，比如写下这本小书，或者和大家交流一下心得。

这本小书是我和 AI 合作完成的，前后花了两个多月的课余时间。因为时间比较仓促，所以篇幅和内容有限，也存在很多不足。其中如有错误或者不当之处，请不吝指出。如果你觉得有收获，可以分享给自己的亲朋好友。如果你有任何疑问，或者想要和我交流，请随时联系。

最后再次提醒，本书所有的推演与分享均基于个人的认知边界与特定时期的观察，切勿将书中的任何推算或历史数据作为现实操盘的机械依据。请时刻对市场保持敬畏，在自己的能力圈内独立思考、审慎决策，走好属于自己的财富与人生之路。

感谢大家支持！

\end{document}
